\begin{chapterbox}
    \chapter{Die Agren, der Schildzwerg und der Feuerdrache (2020)}
    \label{Die Agren, der Schildzwerg und der Feuerdrache (2020)}
    \az{Jahr –921}

    \begin{center}
        einst versteckt unter den Codes 53074 und 54154
    \end{center}
    
    Die Agren Gouri erlebt mit, wie der Schildzwerg Prip vom Feuerdrachen Sagrak angegriffen wird. Sie trifft eine mutige Entscheidung, die sie direkt in den Unterirdischen Krieg zwischen Zwergen, Drachen und Trollen hineinzieht. Und da lauern eine uralte Zwergenfeste, ein Meisterschmied mit seinen mächtigen Schilden, der magische Ort Krahal und sogar der gefährlichste Drache aller Zeiten.
\end{chapterbox}

\section{Gouri, Prip und Sagrak}

\az{Jahr –921}

„DRACHEE!“

Der verzweifelte Schrei hallte durch die Täler des Grauen Gebirges und schreckte allerlei Kleingetier auf, welches schleunigst das Weite suchte, selbst wenn es die Warnung nicht verstehen konnte. Die wenigen Agren, die sich ausserhalb ihrer sicheren Höhlen aufhielten, schreckten auf. Sie rannten allerdings nicht blindlings davon, sondern richteten ihre Augen furchtsam gen Himmel, blickten nach oben, suchten den kleinen schwarzen Punkt, welcher Feuer und Tod ankündigen würde.

Gouri, eine vergleichsweise junge Agren (sie hatte doch erst kürzlich ihren vierzigsten Jahrestag gefeiert), versuchte ebenfalls, die Gefahr im Himmelszelt zu erkennen. Den Kopf in den Nacken gelegt, starrte sie nach oben, direkt in die blendende Sonne. Sie senkte ihren Blick auf die übrigen grünhäutigen Angehörigen des friedlichen Bergvolkes, die ein ganzes Stück von ihr entfernt (leider) die Pflanzenwachsstätten umgegraben hatten. Normalerweise mochte Gouri die Zeit alleine im Freien, wo man in der Mitte einer Wiese sitzen und über den Lauf der Welt nachsinnieren konnte. Nun stand sie aber alleine auf weitem Felde, und der Drache könnte sie von allen Seiten überraschen. Furcht breitete sich in ihrer Magengegend aus.

Da! Ein weiterer Schrei ertönte, diesmal von einem der Agren auf den Pflanzenwachsstätten! Die übrigen Agren rannten vereint los in Richtung einer der rettenden Höhlen, sorgsam darauf achtend, niemanden zurückzulassen. Niemanden ausser Gouri, von der sie nicht einmal wussten, dass sie sich hier draussen aufhielt. Die verängstigten Agren rannten vor Gouris aktueller Position weg. Diese war sich bewusst, was das zu bedeuten hatte: Der Drache befand sich irgendwo hinter ihr!

Innerlich schon mit ihrem Leben abschliessend, drehte sich Gouri um, um einen atemberaubenden Anblick zu erleben. Ein riesiges Reptil, dessen schwarze Schuppen im hellen Sonnenlicht glänzten und glitzerten, brach im Sturzflug durch eine nahe Wolke hindurch. Was zunächst ein kleiner Fleck am Himmel gewesen war, wurde rasant grösser, bis Gouri sogar die gebogenen Zähne im offenen Maul erkennen konnte. Alle Instinkte in ihr schrien sie an, sie solle losrennen, sich vor dieser wandelnden Naturkatastrophe verstecken, irgendetwas tun.

Gouri wusste es besser. Sie blieb zitternd stehen und versuchte, sich zu fokussieren.

In ihrem Kopf bereitete Gouri sich bereits auf den Mentalen Schrei des Drachen vor, welcher jeden Moment ertönen konnte. Sie konzentrierte sich auf ein Kinderlied, welches ihre Mutter ihr jede Nacht vorgesungen hatte, als die Welt noch heil gewesen war – ehe der Meisterschmied Kreatok ein Massaker unter den Drachen angerichtet hatte und der Krieg ausgebrochen war, welcher seitdem das Graue Gebirge plagte.

So sang Gouri innerlich:

\textit{‚Eia, Heja und Naidoa, der Nordwind, der Ostwind und Westwind sind da...’}

Da sandte der angreifende Drache auch schon seinen Mentalen Schrei gegen alle Wesen in seiner Nähe aus:

„MÖRDER! DIEBE! VERRÄTER! VERROTTEN MÖGT IHR, IHR UND ALLE, DIE EUCH ZU BESCHÜTZEN VERSUCHEN!“

\textit{‚Eja, Heja und Naidoa, die Winde, die sausen und brausen durch Tal...’}

Solange Gouri sich auf den Gutenachtreim konzentrierte und so den geistigen Angriff des Drachen verdrängte, gelang es ihr, wenn auch schwankend, auf den Beinen zu bleiben. Vielen anderen Wesen in der Nähe ging es allerdings nicht so. Insektenhirne dachten zu wenig, um gross vom Mentalen Schrei betroffen zu sein, doch Vögel fielen links und recht wie von Steinen getroffen vom Himmel, unfähig, viel mehr zu tun als das Kreischen des Drachen zu erleiden. Dann waren da die Kaninchen, die Rehe, die Füchse, die Wölfe der Gegend, die sich am Boden wälzten und verzweifelt versuchten, etwas zu verscheuchen, was man nicht verscheuchen konnte. Und dann gab es noch den armen gepanzerten Schildzwerg, welchen Gouri in jenem Augenblick entdeckte.

Gouri wusste nicht, ob es daran lag, dass der Zwerg nie gelernt hatte, eine mentale Attacke abzuwehren, oder ob der Drache sich in seinem Schrei ganz besonders auf ihn fokussierte, aber der Zwerg steckte das Ganze erheblich schlechter weg als Gouri. Er war auf die Knie gefallen und hatte sich dann zusammengerollt, unregelmässig atmend, die Zähne zusammengebissen und die Hände auf die Ohren gepresst – doch was konnte das schon ausrichten, wenn der Schrei nicht durch die Ohren in seine Wahrnehmung trat?

„LUMPENGESINDEL! DAS IST FÜR NEHAL! DAS IST FÜR SERAFIM! DAS IST FÜR NERUPAK! SIE ALLE SIND GEFALLEN, EURETWEGEN! IHR BLUT KLEBT AN EUREN DRECKIGEN HÄNDEN UND DUNKLEN SCHILDEN!“

Einen Mentalen Schrei aufrechtzuhalten kostete Unmengen an Energie, lange konnte der Drache dies bestimmt nicht mehr durchstehen. Das war aber auch nicht sein Ziel. Noch immer im freien Fall breitete er seine ledernen schwarzen Flügel aus – der abrupt darin gefangene Wind ruckte schmerzhaft an seinen Gelenken – fing seinen Fall ab und wehte dabei ein ganzes Dutzend kleinerer Bäume um. Schnell jetzt, ehe diese Ratte von einem Wicht sich wieder fassen konnte! Der Boden erzitterte unter seiner Landung. Der glorreiche schwarze Drache – Sagrak war sein Name – richtete seinen Schlangenhals auf und spürte verzückt das vertraute warme Brodeln, wie es von seinem Magen langsam seinen Hals hinaufkletterte und seinen Schlund erreichte. Sagrak blähte seine Nüstern, öffnete seinen Rachen und grinste in Antizipation dessen, was er diesem mickrigen Spornenfrass vor ihm gleich antun würde. Dann spie er einen Strahl puren Drachenfeuers, golden glänzende Flammen, denen nichts und niemand standhalten konnte.

Gouri konnte deutlich die bösartige Freude fühlen, die in Wellen vom Feuerdrachen ausströmte, und sie war sich sicher, dass der Schildzwerg dies auch spürte. Den Mentalen Schrei hatte der Drache – Sagrak? – offenbar schon gemeistert, aber die Fähigkeit, seine Emotionen zu beherrschen und in seinem eigenen Kopf zu verwahren, war ihm noch nicht gegeben. Er musste ein Jungtier sein. Darauf deutete auch seine Grösse hin – selbst der Zwerg reichte ihm bis an die Schulter. Nichtsdestotrotz war auch ein Jungdrache eine nahezu unüberwindbare Gefahr.

Unter enormer Anstrengung klaubte der immer noch am ganzen Leib zitternde Zwerg jetzt einen angeschlagenen Buckelschild hervor und wuchtete ihn in letzter Sekunde zwischen seinen ungeschützten Körper und den Flammenwall, welchen der Drache auf ihn spie. Drachenfeuer war heisser als jedes andere natürlich vorkommende Feuer dieser Erde, sodass auch der beste Schild in diesem Inferno bald vernichtet gewesen wäre. Und der Schild, hinter welchem der Schildzwerg sich verbarg, war wahrlich nicht der beste Schild. So verbrannte das Holz und das Eisen schmolz buchstäblich dahin, während der arme Schildzwerg gezwungen war, unten am Boden zu bleiben und zu hoffen, dass der Schild ihn noch lange genug schützen würde. Lange genug, bis der Drache wieder Luft holen musste. Lange genug, um am Leben zu bleiben. Gut sah es nicht aus. Die Haare an seinen Armen waren schon längstens verglüht, und sein mächtiger brauner Bart würde als nächstes folgen, das wusste er. Gouri sah, wie er die Augen schloss und ein Stossgebet an die Mutter des Steins sandte. Wartete er auf das Ende dieser elenden Qual, welches auch sein Ende sein würde?

Da klappte Sagrak seinen Schlund zu und der Feuerstrahl versiegte. Schwefelgeruch lag in der Luft. Tief durchatmend stand der Drache da, und sowohl Gouri als auch der Schildzwerg nahmen seine Erschöpfung und Verwirrung wahr – er hatte eindeutig seine Feuerkraft überschätzt. Das sollte aber nicht heissen, dass er des Zwerges Lebenslicht nicht immer noch durch einen mächtigen Hieb mit einer seiner Pranken einfach auslöschen könnte.

Gouri beobachtete gebannt und immer noch zittrig die Szene, die sich vor ihr abspielte. Sie war sich uneinig, was sie nun tun sollte. Jetzt war ihre Gelegenheit, sich umzudrehen und in die sichere Haushöhle zurückzukehren. Zurück zu ihren Eltern, die sich bestimmt schon riesige Sorgen machten. Die ihr wahrscheinlich verbieten würden, innerhalb der nächsten drei Sonnenzirkel das Tageslicht wieder zu betreten. Gouri seufzte.

Dieser Schildzwerg... er schuldete ihr nichts, denn es waren die Zwerge gewesen, die ihr Bündnis mit den Drachen gebrochen und so nebst der Zerstörung vieler ihrer so prächtigen Bauten auch Krieg und Unheil über das einst so friedliche Graue Gebirge gebracht hatten. Eigentlich sollte Gouri sogar froh sein, dass es bald einen Zwerg weniger auf dieser Erde gab. Aber sie konnte nicht. Der Mut dieses Kriegers faszinierte sie: Er hatte einen Drachen auf sich gehetzt, einen Mentalen Schrei erlitten, und er lebte immer noch. Er hatte seine Waffen verloren, sein Schild war verkohlt und sein eleganter Bart hatte sich im Flammeninferno von eben gerade verabschiedet. Dennoch stand er vorsichtig auf, aufrecht dem sich immer noch von der Anstrengung erholenden Drachen entgegenstellend, und blickte ihm entschlossen und ungerührt ins echsenartige Angesicht.

In ihrer Kindheit hatte Gouri stets fasziniert den Legenden von Zwergen und Drachen, von mächtigen Schilden und den Landen jenseits des Grauen Gebirges gelauscht. Und die Heroen in diesen Geschichten liefen nicht von einem unfairen Kampf davon, sondern setzten sich für Gerechtigkeit ein. Und dieser Kampf zwischen Schildzwerg und Drache war nicht gerecht.

In diesem Moment entschloss Gouri sich: Sie würde nicht wegrennen. Sie würde wenigstens diesen einen Streit zwischen Zwerg und Drache zu schlichten versuchen! Das fühlte sich richtig an! Und so richtete Gouri sich auf, ihr verfilztes Agrenhaar wehend im heissen Wind, und war schon kurz davor, zwischen den Schildzwerg und den Feuerdrachen zu springen, um diesem Unsinn Einhalt zu gebieten...

...da fasste sie sich wieder. Einfach so zwischen einen wütenden Drachen (insbesondere ein Jungtier) und einen Schildzwerg zu springen, war kein gescheiter Weg, sein Leben zu verlieren. Gab es vielleicht andere Wege, wie sie dem ungeschützten Schildzwerg helfen konnte?

So jung, wie Sagrak, der Drache, war, wollten seine Eltern bestimmt nicht, dass er sich alleine auf Zwergenjagd befand. Auf einem aufmüpfigen Ausflug war er also! Wahrscheinlich kannte er sich nicht im Geringsten mit der realen Welt aus. Er hatte wohl davon geträumt, bei seiner kleinen Exkursion eine Horde böser Zwerglein zu finden und problemlos zu flambieren, „ihrem gerechten Ende zuzuführen“, wie die Drachen es in ihren Gutenachtgeschichten ihren Kindern beibrachten. Und nun wunderte er sich, wie es sein konnte, dass er den Schrei genutzt und seinen Feueratem verbraucht hatte, während dieses dreckige Wichtlein vor ihm noch lebte. Viel brauchte es nicht mehr, damit Sagrak verunsichert das Weite ziehen müsste. Allerdings nichts, was die schwache Gouri selbst tun könnte. Aber vielleicht jemand anderes?

Gouri blickte wild um sich. Die meisten Wesen in ihrer Nähe waren geflohen, doch da! Dort im Felsen hinter Sagrak konnte sie einen Eingang zu einem Arpachenbau erkennen!

Die Arpachen, riesige Insektoide, hatten mit ihren mächtigen Kiefern einst Gänge durch das Graue Gebirge gegraben und wie die Zwerge ein System aus unterirdischen Gängen gebaut – wenn auch weit weniger elegant. Im Gegensatz zu den inzwischen aus taktischen Gründen zugeschütteten Eingängen zum Zwergenreich Cavern waren die Höhlen, die die Gänge zu den Arpachenbauten markierten, weder getarnt noch verschlossen, und so nutzen verschiedenste Spezies sie gerne als Aufenthaltsort.

Allen voran natürlich die Agren, die Gouri in so einer Situation nicht gross helfen könnten.

Ein Höhlenwicht wäre auch nicht wirklich hilfreich; diese schleimigen Gestalten waren kaum stärker als ein Agren und griffen bloss an, wenn ihr Gegner ihnen den Rücken zuwandte.

Eine Sporne wäre schon etwas besser; diese Spinnenwesen konnten zwar nichts gegen einen Drachen ausrichten, würden aber, wenn provoziert, wütend aus ihrer Höhle schiessen und durch aufgestellte Vorderbeine versuchen, ein möglich einschüchterndes Bild abzugeben, was vielleicht schon genug war, um Sagrak zu vertreiben.

Ein Hornbär wäre noch effektiver; diese gehörnten Riesenbären konnten es tatsächlich mit einem jungen Drachen aufnehmen, wenn sie aus ihrem Winterschlaf geweckt wurden.

Der Hauptgewinn wäre aber eine Trollhöhle. Die Trolle waren mit Abstand die stärksten flugunfähigen Wesen im Grauen Gebirge und eine Horde von ihnen hatte sogar schon einmal einen ausgewachsenen Drachen erlegt. Sie waren nicht mit hoher Intelligenz gesegnet – den besiegten Drachen hatten sie verspeist, ohne das Fleisch auch nur anzubraten zu versuchen, und diejenigen Trolle, die nicht an die Lebensmittelvergiftung ihr Leben gelassen hatten, waren kurz daraufhin von den übrigen Drachen dabei erwischt worden, wie sie triumphierend aus den skeletalen Überresten des Gefallenen eine Hütte zu bauten versuchten – aber man brauchte keine hohe Intelligenz, um einen kleinen Drache wie Sagrak zu verschüchtern.

Und so klaubte Gouri einen scharfen Stein aus dem felsigen Boden, wog die Distanz zur mysteriösen Höhle kurz ab und warf ihn dann energisch in die Richtung des Höhleneingangs.

Sie verfehlte um Längen und traf Sagrak am Kopf.

Rückblickend war das vielleicht sogar etwas Gutes, da es dem Schildzwerg einige wertvolle Sekunden schenkte. Aber in dem Moment, als Gouri erblickte, wie ihr Stein den Drachen an der Schläfe traf und dieser blitzschnell den Kopf drehte und sie aus seinen glühenden blutroten Augen direkt anstarrte – in diesem Moment hüpfte Gouris Herz in ihre Mooshose, und sie konnte nichts mehr, als „Spornendreck“ zu flüstern und innerlich mit ihrem Leben abzuschliessen. Dann fasste sie sich wieder, und als der Feuerdrache langsam auf sie zuschritt, den Kopf auf dem langen Hals weit erhoben (sie wusste, dass er sich so auf einen erneuten Feuerstrahl vorbereitete, und dieser wäre ihr Ende und das des Schildzwergs gewesen), der lange stachelbesetzte Schwanz zuckend vor Vorfreude, warf Gouri voller Adrenalin einen zweiten Stein.

Der spitze Gesteinssplitter drehte sich mehrmals um die eigene Achse, beschrieb eine elegante Flugparabel über Sagraks Kopf hinweg und verschwand im Dunkel der Höhle, wo ein leises Klack anzeigte, dass er ein Ziel gefunden hatte.

Ein tiefes, wütendes Brüllen erklang. Wie viel Glück konnte eine Agren haben? Es handelte sich tatsächlich eine Trollhöhle!

Sagraks Augen verengten sich, und er wandte sich ebenso blitzschnell, wie er seine Aufmerksamkeit auf Gouri gerichtet hatte, wieder von ihr ab. Gouri sah in ihm nicht mehr das feuerspeiende, bösartige Monster, dass sie noch eine Sekunde zuvor in ihm gesehen hatte. Stattdessen war er bloss ein armes Jungtier, hineingeführt in einen Konflikt, welchen es nicht begonnen hatte und welchen es nicht beenden konnte. Furcht stand in Sagraks Blick und seine Furcht brach aus seinem Kopf heraus, erfüllte Gouri und den Schildzwerg.

Dann trat der Troll aus der Höhle.

Es war ein furchterregendes Monstrum, Beine so dick wie Eichenstämme, Muskeln wie ein Bär und riesige gewundene Hörner, für die es keinen Vergleich im Tierreich gab – nicht mal ein Hornbärenhorn konnte mit der Stabilität und Tödlichkeit eines Trollhorns auch nur annähernd mithalten. Jetzt brüllte der Troll, sich immer noch das Auge reibend, welches Gouri anscheinend getroffen hatte: „WER WAGEN WERFEN DIESEN STEIN HIER?!“

Trolle brauchten nicht viele Worte, damit man sie verstand. Der breite Troll blickte sich wütend um und fokussierte seinen Blick natürlich auf den grössten potentiellen Gegner im Raum, Sagrak.

Ungeachtet der Tatsache, dass Sagrak mit seinen Tatzen wahrscheinlich nicht mal einen grossen Schild greifen konnte, also definitiv nicht den Steinbrocken geworfen hatte, stapfte der Troll furchtlos auf den Drachen zu und brüllte: „DIESER TROLL HIER KNACKEN DIESEM DRACHEN DORT SEIN HALS!“

Die beiden Riesenwesen machten Augenkontakt. Weisse, pupillenlose Augen trafen auf loderndes Rot. Dann war der Troll bei Sagrak angekommen. Nun spie Sagrak Feuer, doch der Troll war schneller und drückte ihm die Schnauze zu. Ein klein wenig Glut floss dem prustenden Drachen aus der Nase, und der Troll brüllte auf vor Schmerz, als es über seine Hand floss und seine Haut versengte. Er hielt jedoch weiterhin fest. Sagrak versuchte verzweifelt, sich loszustrampeln. Er wandte und drehte sich, und entschlüpfte dem Griff des Trolles. Seine Flügel bereits entknittert, sprang er hoch in die Luft, um entrüstet davonzufliegen.

Gouri war damit zufrieden. Sobald der Drache in den fliegenden Zustand gekommen war, würde ihm keine Gefahr mehr durch den Troll drohen. Sagrak würde verstört das Weite suchen, und der Troll würde ihm folgen, von der einfältigen Hoffnung erfüllt, ihn immer noch zu erwischen. Sie und der Schildzwerg wären in Sicherheit. Ende gut, alles gut.

Gouri stolperte den Hang hinunter auf den sich immer noch hinter Schildüberresten duckenden Zwerg zu und näherte sich ihm vorsichtig. Seine Augen weiteten sich, als er sie als Angehörige der Agren erkannte, und sein Blick schien sie anzuflehen, die Dummheiten sein zu lassen und sich nicht seinetwegen in Gefahr zu begeben. Gouri ignorierte sie und wollte ihm schon aufhelfen...

...da packte der Troll den wegfliegenden Sagrak am dornigen Schwanz und zog diesen zurück auf den Erdboden, wo er ungeschickt aufklatschte. Die Erde erbebte, sodass selbst Gouri Mühe hatte, auf ihren stämmigen Agrenbeinen stehen zu bleiben. Der Troll überquerte den Weg zu Sagraks Brustkorb mit einem gewaltigen Sprung und landete auf ihm, seine massigen Beine links und rechts auf die zerschmetterten Flügel der Echse stellend. Sagrak schrie auf. Sein Schrei ging allen Anwesenden durch Mark und Bein. Dann schnappte er mit seinen scharfen Zähnen nach der Nase des Trolles. Dieser lachte nur auf, klemmte Sagraks Schnauze erneut fest. Sagrak stach mit seinem Schwanz nach dem Rücken des Trolls. Dieser grunzte nur amüsiert und drückte noch fester zu.

Die kämpfenden Giganten waren Gouri definitiv zu nahe – wenn Sagrak eine ungeschickte Bewegung mit seinem Dornenschwanz machen würde, dann wären sowohl sie als auch der Schildzwerg Spiesschen. Ein Stossgebet sendend an die Geister des Waldes, des Steins und der Lüfte, näherte sich Gouri dem angekokelten Schildzwerg, welcher sich inzwischen erhoben hatte.

„Folge mir!“, flüsterte sie eindringlich. Der Zwerg tat wie geheissen und humpelte ihr nach, weg von den kämpfenden Biestern. Sie mussten zu einem sicheren Agrenbau, wo den Schildzwerg hoffentlich Verpflegung und Arzneien erwarten würde, und zwar schnell!

Als sie glaubte, genug weit weg von Troll und Drache zu sein, warf Gouri einen Blick zurück.

Ihr gefiel nicht, was sie erblickte.

Sagrak wand sich ein letztes Mal vergeblich. Verzweifelt sandte er einen leisen Mentalen Schrei aus, welcher nicht einmal den angeschlagenen Schildzwerg von den Füssen holte. Dann blubberte er einen letzten Rest Lava aus seinem Magen heraus und versuchte, diesen dem Troll anzuspucken. Nichts half ihm aus seiner misslichen Lage. Sein Leben zog vor seinem inneren Auge vorbei, und so offen, wie sein Geist momentan war, zog sein Leben auch vor den Augen von Agren, Schildzwerg und Troll vorbei. Kurze Eindrücke von verschiedensten Momenten schossen auf sie ein: Der erste Riss in Sagraks Eierschale, durch den das helle Licht der Aussenwelt schien; seine erste Jagd; seine erste Flamme, wie sie einen ganzen Baum mit nur einem Niesser in Brand setzte; ein Bad im See am Grunde der Krahal-Schlucht mit einer Drächin, die er sehr mochte; die reine Wut, als er von Kreatoks Verrat gehört hatte; die Genugtuung, mit welcher er Dutzende Zwerge durch die Winterburg jagte; die Furcht, als direkt neben ihm sein bester Freund mit einem riesigen Stein aus einem Wurfgeschoss der Riesen vom Himmel geschlagen wurde; die Rachegefühle, die er verspürte, als er die Zera tief im Lavafluss Enran begrub, gemeinsam mit den zwei Riesen, die das Katapult bedient hatten... dann wurde es still.

Äusserlich hatte sich am Drachen nichts verändert, doch das Feuer in seinen Augen war erloschen, und er wehrte sich nicht mehr. Sein Körper wurde schlaff und seine Klauen bohrten sich nicht mehr länger in die grobe Haut des Trolls, sondern fielen kraftlos auf die Seite.

Die letzten Eindrücke, die die Anwesenden aus Sagraks Geist empfangen hatten, waren ein Bild und ein einziges Wort gewesen.

Das letzte Bild, welches dem Drachen durch den Kopf geschossen war, war ein riesiger schwarzer Baum aus Stein, über und über mit Edelsteinen verziert, in einer Art Höhle, in die wohl die ganze Winterburg doppelt und dreifach gepasst hätte. Der Steinbaum war umgeben von einem sanften blauen Leuchten, und dutzende von Drachen sassen auf seinen Ästen und unterhielten sich. Weiter unten sah man kleine Kreaturen, die mit ihren Hornklauen die Rücken und Flügel einiger seufzender Drachen massierten. Eine friedliche Stimmung herrschte.

Das letzte Wort, welches der Drache gedachte hatte, war ‚Krahal’.

Gouri hatte den Verdacht, dass Sagraks Geist nicht mehr in diesem Körper weilte, als der Troll triumphierend den Hals des leblosen Drachen schnappte und mit einem breiten Grinsen einmal um die eigene Achse drehte, sodass die Wirbel nur so knackten. Ihr wurde übel und sie musste sich von der Szenerie abwenden. Ein Schleifgeräusch verriet ihr, dass der Troll, dessen Kampfeslust soeben befriedigt worden war, gerade stolz seinen Fang in die Höhle zurückzerrte, um vor seiner Familie damit anzugeben. Mögen sie doch dran ersticken!

Gouri war selbst über ihre Wut auf den Troll überrascht, die sie verspürte. Immerhin hatte der Troll ihr Leben und das des Zwerges gerettet. Aber er hatte auch ein Leben beendet, und er freute sich daran. Das hätte der Drache zwar auch getan, aber dieser hatte immerhin eine gute Begründung für seine Taten. Nun gut, gut war seine Begründung ja nicht, aber wenigstens verständlich. War das Verhalten des Trolles etwa nicht verständlich? Verwirrt liess Gouri sich zu Boden fallen, während ihre Gedanken weiter kreisten.

Jetzt war der Schildzwerg an ihrer Seite. Seinen gepolsterten Helm mit den vielen Runenverzierungen hatte er abgesetzt. Die Überreste seiner Rüstung, sein Bart und sein Haupthaar waren grösstenteils weggebrannt – unter seiner Rüstung erkannte Gouri eine isolierende Schicht schwarzer Stoffstreifen, welche er sich umgebunden hatte – doch strahlte er immer noch etwas Erhabenes aus, wie er vor ihr stand und sprach: „Der Mutter des Steins sei Dank. Nein, Euch sei Dank! Was Ihr getan habt... ich dachte, die Agren mischen sich nicht ein in... egal, ich will nur ausdrücken... wie kann ich das nur je wiedergutmachen?“

Gouri war immer noch schockiert von der Gewalt, die sie gerade erlebt hatte. Sie hatte die Geschichten gehört und einige Drachenangriffe aus der Ferne miterlebt, doch dies hier war ein Erlebnis ganz neuen Ausmasses gewesen.

„Diese Trolle sind eine elende Plage!“, stiess sie hervor. Das machte keinen Sinn, der Troll hatte soeben ihr und des Zwerges Leben bewahrt. Sie sollte nicht wütend auf ihn sein. Aber die sinnlose Grausamkeit, die das Wesen nonchalant an den Tag gelegt hatte, weckte überraschend viel mehr Emotion in ihr als die Grausamkeit, die die Zwerge und Drachen einander aufgrund ihrer langjährige Feindschaft zeigten. Nicht, dass der Krieg sinnvoll gewesen wäre.

Der Zwerg erwiderte: „Ihr meint, das Graue Gebirge hätte eine Trollplage? Ich komme gerade von einem Treffen aus dem Trollland in Norden. Die dort lebenden Hünen überragen den grössten Schildzwerg um über zwei Köpfe und leben dennoch im wahrsten Sinne des Wortes in den Baumwipfeln, weil auf dem Boden zu viele Angriffe durch Trolle drohen. Das ist doch verrückt! Das nenne ich eine wahre Trollplage!“

Er klang beinahe enerviert. Dann fiel ihm auf, dass so eine Reaktion nicht förderlich war, und fügte hinzu: „Verzeiht bitte. Es war ein... aufregender Tag.“

Gouri schwieg. Sie wusste nicht, was zu sagen war. Sie wusste nicht, wie es jetzt weitergehen sollte. Sie hatte stets davon geträumt, selbst Teil einer Sage zu werden und einen jungen Zwergenprinzen zu treffen. Aber die Realität war auf einmal überwältigender als ihre Fantasie.

„Ich bin Gouri“, sprach die Agren schlicht.

„Prip, Sohn des Aigar, Bote des Eisernen Stuhls von Cavern und Hüter des Casamatucs der Drei Wasser, zu Euren Diensten“, erwiderte der Schildzwerg mit einer angedeuteten Verbeugung.

Gouri wünschte sich, sie hätte auch einige Titel zu nennen. Sie streckte ihm ihre grüne Hand zur Begrüssung hin und Prip schüttelte sie mit seinen Eisernen Handschuhen. Gouri schrie auf, denn die Handschuhe waren immer noch heiss vom Odem des Feuerdrachen (offenbar waren sie innerlich isoliert, sodass Prip die Wärme nicht verspürt hatte). Prip entschuldigte sich vielmals. Beide tauschten ein nervöses Lachen aus. Dann besann Gouri sich: Ihr ging es immer noch prima, doch Prip war, auch wenn er es sich nicht anmerken lassen wollte, verwundet und verbrannt, und brauchte dringend einen Heiler. Sie mussten zurück zu den anderen Agren.

„Prip, Dir... Euch geht es nicht gut. Ich kenne den besten Heiler des Grauen Gebirges und kann Euch zu ihm führen. Folgt mir!“

Prip blickte sie kurz stolz an, als wollte er sie herausfordern, seinen gesundheitlichen Zustand nicht zu kritisieren. Dann sprach er: „So gerne ich auch die Dienste dieses Heilers in Anspruch nehmen würde, nehme ich doch an, dass er den Schildzwergen ein wenig zurückhaltender gegenübersteht als Ihr, und ich will die Gastfreundschaft der Agren nicht überstrapazieren. Zudem muss ich dringend zu der Zwergenfeste „Die Drei Wasser“, dem Fürsten von Cavern Bericht erstatten gehen.“

Gouri prustete beinahe los. Prip würde in seinem jetzigen Zustand doch keine drei Baumlängen weit kommen. Allerdings hatte er recht – die übrigen Agren würden sich nur ungern um einen Schildzwerg kümmern wollen, erst recht nicht, wenn dieser gerade erst eine Auseinandersetzung mit einem Drachen gehabt hatte. Auch wenn es die Drachen waren, die zwischendurch den einen oder anderen unvorsichtigen Agren brateten, so waren es in ihren Augen die Zwerge, die die Angriffe dieser Tiere provoziert hatten, und so waren Zwerge keine gern gesehenen Gäste.

Gorui fasste einen Entschluss: „In diesem Fall werde ich Euch begleiten. Unterwegs stossen wir bestimmt auf einige Büschel Wolfskraut, dann kann ich Euch immerhin mit einer behelfsmässigen Salbe aushelfen.“

„Wie tief soll ich denn noch in Eurer Schuld stehen?“, ächzte Prip.

Gouri bemerkte, dass er nicht protestierte. Sie verspürte einen Anflug kindlicher Freude. Sie würde diesen Zwerg zu den Hallen der Zwergenkönige begleiten! Sie würde auf ein Abenteuer gehen!

Erst als die beiden die Alte Zwergenstrasse erreichten und eine weitere Höhle passierten, fiel Gouri ein, dass die anderen Agren wohl gerade umkamen vor Sorge um sie.






\newpage
\section{Der Sturm auf die Drei Wasser}

„Obacht! Da vorne ist etwas!“, zischte Gouri zu Prip. Prip verharrte, wo er war, und packte die verkohlten Überreste seines Schildes fester. Schildzwerge, dachte Gouri belustigt, immer auf ihre Waffen fixiert. Dieses verbrannte Stück Holz war die einzige Waffe, die ihm geblieben war, und nun hielt er daran fest, als hinge sein Leben daran.

„Ein Lebewesen?“, wisperte Prip angespannt. Zwerge wie er konnten im Dunkel der Nacht keine fünf Hasensprünge weit blicken, wenn sie nicht eine ihrer genialen Kerzenbrillen trugen. Gouri richtete ihre erheblich besser an die Dunkelheit angepassten Agrenaugen nach vorne, in die Richtung, aus der das Geräusch gekommen war. Im dichten Nebel konnte sie aber auch nicht viel mehr einen grossen Schemen ausmachen, der sich kaum bewegte – natürlich ausgerechnet auf der anderen Seite der Hängerücke, die sie und Prip überqueren mussten.

„Es bewegt sich kaum, aber es bewegt sich. Das wird ein Lebewesen sein.“, flüsterte Gouri zurück.

„Was kannst du erkennen?“

„Es ragt aus dem Nebel heraus. Es ist gross.“

„Ein Troll bei einem Mitternachtsspaziergang vielleicht? Oder ein Drache?“

„Auf jeden Fall ist es kein gutes Zeichen.“

Gouri bedachte ihre Optionen. Sie konnten versuchen, die Brücke zu überqueren, aber das würde das Wesen auf jeden Fall auf sie aufmerksam machen. Und zu warten, könnte sinnlos sein, falls sich der Schemen letzten Endes nur als seltsame Felskonstruktion herausstellen sollte. Mühsam unterdrückte sie ein Gähnen. Agren brauchen ihre dreizehn Stunden Schlaf, um fit und munter zu sein, und Gouri war momentan weder das eine noch das andere.

„Gibt es andere Wege zu den Drei Wassern?“, fragte sie Prip.

„Es gibt einen Pfad der Korn-Schlucht entlang bis zu einem Felsvorsprung, von wo aus wir zum Noswald springen können. Aber dieser Umweg würde uns das einen ganzen Tag kosten. Mal ganz abgesehen davon, dass ich in meinem jetzigen Zustand den Sprung kaum schaffen werde.“

Gouri war überrascht, sonst versuchte Prip doch immer, so zu tun, als wäre er nicht am ganzen Körper versengt und verkratzt. Offenbar war es ihn wirklich wichtig, so schnell wie möglich die Drei Wasser zu erreichen. Wie um ihre Gedanken zu bestätigen, betonte Prip noch einmal: „Es ist wirklich dringend. Wir müssen diese Brücke jetzt überqueren.“

Gouri seufzte. Eine gute Mütze Schlaf kam wohl nicht in Frage. Dabei sahen diese Moosbüschel selbst im Dunkeln unglaublich einladend aus...

Als hätte Prip ihre Gedanken gelesen, fügte er an: „Von der anderen Seite der Brücke ist es ein Katzensprung bis zu den Drei Wassern. Dort wirst du die Gastfreundschaft des Zwergenvolkes erleben dürfen, erlesene Speisen, Federbette, warmes Kaminfeuer!“

Gouri antwortete nicht. Prip liess seine Schultern hängen und fuhr fort: „Du hast ja recht, du musst die Brücke nicht überqueren, wenn du nicht willst. Ich danke dir vielmals für alles, was du bislang schon für mich getan hast, ohne dich und deine Heilkünste wäre ich wohl keinen Hasensprung weit gekommen. Ich werde dich nicht vergessen und eines Tages werde ich zurückkommen und meine Schuld...“

Gouri atmete tief durch und stiess dann hervor: „Jetzt schweig schon stille, ich komme ja weiterhin mit.“

Prip hatte klargemacht, dass er auf jeden Fall versuchen würde, diese Brücke zu überqueren, und Gouri wollte den schwachen und angeschlagenen Schildzwerg nicht alleine durch diese düstere Nacht wandern lassen. Und Prips Worten nach lagen die Drei Wasser und weiche Zwergenbette nicht weit entfernt hinter dieser Schlucht. Ganz abgesehen davon wollte Gouri auch nicht alleine im Dunkeln zurück zu ihrer Agrenhöhle wandern, und darüber nachdenken, was für Sorgen sich ihre Familie gerade machen musste, wollte sie auch nicht. Von den Drei Wassern konnte sie bestimmt einen Falken mit einer Botschaft an die Agrengesellschaft schicken. Falls den Zwergen ihre Wünsche überhaupt wichtig waren. Aber darum konnte sie sich später kümmern, zunächst einmal hiess es, diese verflixte Hängebrücke zu überqueren.

Wenn sie nur nicht so viel Angst davor hätte, in die Tiefe zu stürzen.

Prip beobachtete Gouri besorgt, während er einige Schritte auf die Brücke zuging. Gouri folgte ihm nur langsam. Prip drehte sich zurück zu ihr und fragte: „Ist etwas nicht in Ordnung? Du musst mich wirklich nicht über diese Brücke begleiten.“

Gouri atmete etwas schneller. Grosse Höhen waren noch nie ihr Ding gewesen, und Hängebrücken erst jetzt nicht, aber hier im Dunkeln, da war noch etwas anderes. Diese grosse Gestalt im Nebel auf der anderen Seite der Brücke weckte in ihr eine zusätzliche Furcht, die Prip einfach nicht gleich zu spüren schien. Oder es gut versteckte. Gouri wollte diese Brücke einfach nicht überqueren. Sie murmelte etwas über Höhen und das Dunkel und Prip schien kurz abzuwägen, ob er sich trotzdem einfach von ihr verabschieden sollte und den Rest des Weges zu den Drei Wassern alleine zurücklegen. Dann sprach er:

„Wie wär’s damit: Wir gehen einige Schritte auf die Brücke zu, und schauen, ob sich die Gestalt bewegt. Wenn ja, dann rennen wir so rasch wie möglich zurück in den Nebel, wo uns niemand erwischen wird. Wenn nicht, dann können wir weitergehen.“

Gouri nickte stumm. Dann fügte sie leise hinzu: „Können wir jeweils alle paar Schritte warten, ob dieses Wesen, falls es denn eines ist, sich bewegt?“

Prip ergänzte sogar: „Wir können jederzeit warten, bis du dich sicher fühlst.“

Gouri war sich nicht sicher, wie sehr das stimmte. Der Zwerg kam ihr schon sehr ungeduldig vor.

Und so betraten die beiden die Brücke, der verletzte Schildzwerg stramm voranschreitend, die unverletzte Agren zaghaft hinterherschleichend.

Drei Holzplanken weiter blieb Gouri stehen, das dicke Brückentau fest im Griff, und lauschte. Sie hörte das Knarren der Brückenbretter und in weiter Entfernung den Schrei einer Horneule. Die Gestalt am anderen Ende der Brücke war weiterhin nur als undeutlicher Schemen zu erkennen, und wenn Gouris Augen sie in diesem angespannten Moment nicht täuschten, wurde der Nebel sogar noch dichter. Ihre unnatürliche Furcht vor diesem Wesen schwellte wieder an, und sie zwang sich, ruhig und langsam ein- und auszuatmen.

Prip, der schon einige Schritte weitergelaufen war, drehte sich um und sah Gouri dastehen. Er kehrte zurück zu ihr. „Die Agren, die mir das Leben gerettet hat, indem sie einen riesigen Troll erlegt hat, fürchtet sich vor ein wenig Höhe? Das ist echtes Zwergenhandwerk, diese Brücke wird in einem Jahrhundert noch stehen! Und im Dunkeln sieht man nicht einmal, wie tief es ist...“

Prips Stimme wurde leiser und setzte aus. Er wusste nichts mehr zu sagen, und wollte weiterziehen.

Gouri riss sich zusammen und setzte ein grimmiges Lächeln auf. Dann stapfte sie weiter, einen Schritt nach dem anderen, links, rechts, links, rechts. Ihre Hand huschte am Brückentau entlang, es immer noch als Stütze nutzend, aber nicht mehr festklammernd. Dann blieb sie erneut stehen. Einen Drittel des Weges über die Brücke hatten die beiden jetzt zurückgelegt, und der Nebel auf der anderen Seite der Brücke verdeckte den Schemen jetzt vollständig. Es bewegte sich nichts. Gouri glaubte, im Wind vereinzelte Stimmen zu hören. Ob diese aus dem unterirdischen Zwergenreich stammten oder nur die durch das Heulen des Windes entstehende „Musik der Toten“ war, vermochte sie nicht zu sagen.

„Was siehst du?“, flüsterte Prip angespannt.

„Nichts“, antwortete Gouri niedergeschlagen, „der Nebel ist zu dicht. Aber es bewegt sich wenigstens nicht.“

„Kannst du weitergehen?“

„Nur... gib mir noch einen Augenblick.“

Prip beobachtete Gouri wieder etwas besorgt, als diese in die Knie ging und tief durchatmete. Das Schwanken der Brücke im Wind wurde nicht weniger, das Holz war alt und morsch, und noch wäre es einfacher, sich umzudrehen und zurückzukehren, einen anderen Weg zu gehen. Aber sie mussten weiter.

Gouri richtete sich wieder auf und ging strammen Schrittes weiter. Sie versuchte, das Knirschen und Quietschen der Brückenbretter nicht zu beachten, und zog ihr Lauftempo noch an. Die Mitte der Brücke überquerte sie schnell, es gab nun kein rasches Zurück mehr. Und der Nebel war immer noch viel zu dicht, um die dunkle Gestalt zu erkennen. Bewegte sich da vorne nicht etwas? Gouri hörte nichts mehr als ihre Schritte auf den Holzplanken und ihren eigenen Atem, viel zu laut in der sonst so stillen Nacht, und sie blieb wieder stehen.

Prip rammte sie von hinten, zum Glück schmerzlos, aber nicht ohne ein für Gouri ohrenbetäubend wirkendes Scheppern seiner Rüstung auszulösen. Sie schnappte sich das Brückentau wieder und warf einen Blick in die Tiefen der Schlucht. Dafür schalt sie sich gleich wieder, denn es vervielfachte ihre Panik nur. Prip fluchte in der Zunge der Zwerge.

Beide blieben wachsam stehen und blickten auf die Nebelschwaden vor ihnen. Dann schauderte Gouri und blickte wieder zurück, in die Richtung, aus der sie soeben gekommen waren. Hatte sie nicht soeben gerade ein Geräusch von da gehört? Keine Seite der Brücke war jetzt definitiv sicher.

„Komm, weiter!“, drängte sie Prip, welcher sie auf die Füsse zog. Leise, leise schlichen Zwerg und Agren über die alte Zwergenbrücke.

Der Mond schien hell und beleuchtete die Szenerie in einem Spiel aus Licht und Schatten. Gouri schaute in den Himmel – nur weg vom Boden – und mochte die Ästhetik eigentlich noch. Wenn die Stimmung nicht so gruselig wäre, so würde man diesen nebligen Himmel und die ganze Lage bestimmt als schön bezeichnen. Schon bildete sich in ihrem Kopf die Anfänge eines Agrengedichts, und sie hielt inne. Diesmal sah Prip es kommen und hielt rechtzeitig an, denn der Nebel lichtete sich.

Der Nebel lichtete sich und gab den Blick frei auf einen riesigen Troll, welcher am Ende der Brücke sass und vor sich hin schnarchte.

Gouri zuckte zusammen. Sie hatte das Bild noch deutlich im Kopf, wie der namenlose Troll von vorhin den Drachen Sagrak vom Himmel zog und ihm ein rasches, gnadenloses Ende bereitete. Sie versuchte, gegen den Instinkt anzukämpfen, auf den Fersen umzudrehen und davonzurennen.

Prip fluchte erneut unter seinem Atem, als auch er den Troll am Brückenende entdeckte. Dann lief er an Gouri vorbei und mit den entschuldigenden Worten „Ich muss zu den Drei Wassern, so schnell wie möglich“ rannte er los, scheppernd und klimpernd, direkt auf den Troll zu.

Gouri schalt ihn innerlich, dann fiel ihr mit Schrecken auf, was der Schildzwerg schon viel früher erkannt haben musste: Der Troll regte sich. Sie hatten ihn bereits aufgeweckt, und nun galt es, an dem Ungetüm vorbeizukommen, ehe es sich vollends auf sie fokussieren konnte.

Prip hatte schon die halbe Distanz zurückgelegt, und Gouri überlegte kurz, ihn alleine weiterziehen zu lassen. Doch bei seinem verletzten Tempo hätte sie ihn im Nu eingeholt.

Gouri rannte los, versuchte, so gut wie möglich einfach nur nach vorne zu blicken. Prip humpelte voran, eine Hand an seinem Schild, die andere am Brückentau. Gouri war noch drei Schritte von ihm entfernt, da traf sie eine heftige Windbö, die sie zusammenzucken liess. Ihr Blick liess von Prip ab, welcher sich jetzt wieder weiter von ihr entfernte und nur noch einige Hasensprünge vom Ende der Brücke entfernt war, zum Troll. Sie zuckte zusammen.

Der Troll war vollends erwacht!

Dieser hier war etwas gedrungener als das Exemplar, welches sie heute bereits getroffen hatte, doch deswegen kein weniger ernst zu nehmender Gegner. Er räkelte sich und gähnte, seine gelblichen Augen mickrig aus seinem hervorspringenden Schädel starrend. Gouri konnte nicht erkennen, ob er sie anblickte. Nun hielt auch Prip inne, und rührte sich nicht mehr. Noch bestand die Chance, dass der Troll sich einfach wieder zu Boden sinken lassen würde.

Prip drehte sich zu Gouri um und gestikulierte, sich nicht zu bewegen.

Der Troll liess sich nicht wieder in sein Grasbett sinken, im Gegenteil, er streckte sich zu seiner vollen Grösse und packte das Brückengeländer mit seinen überlangen Armen. Langsam kam sein mächtiger Bauch auf den Rand der Brücke hinzu. Sein Kiefer knackte zweimal. Hatte er soeben sein Essen erspäht? Würde die Brücke seinem unglaublichen Gewicht standhalten können?

Gouri gestikulierte zurück an Prip, er solle Fersengold geben. Sah er sie überhaupt in der Dunkelheit der Nacht? Prip blieb ruhig und musterte das riesige Vieh, welches ihm den Weg zu den Drei Wassern versperrte.

Gouri fluchte. Verdammte Zwergenfeste, warum musste Prip auch so zwingend dorthin wollen?

Ihr hüpfte das Herz noch mehr in die Hose, als sie erkannte, wie Prip den Griff seines Schildüberrestes etwas stärker packte und leicht anhob. Er bereitete sich auf einen Angriff vor. Dieser Idiot! Dieser sture absolute Idiot! Noch ehe Gouri ein passenderes Schimpfwort einfiel, setzte sich Prip auch schon in Bewegung.

Prip rannte humpelnd und laut scheppernd über die wackelige Hängebrücke auf den Troll zu. Ohne zu überlegen, liess auch Gouri das Brückenseil los und setzte ihm nach, doch Prip legte ein für seinen Zustand erstaunliches Tempo an den Tag und Gouri war immer noch wackelig unterwegs, weswegen sich ihr Abstand nur vergrösserte.

Prip legte die letzten paar Schrittlängen in einem mächtigen Satz zurück und rammte den Troll mit den Überresten seines Schildes in den Bauch.

Prip prallte vom Trollbauch glatt ab und konnte sich gerade noch ein Seil der Hängebrücke fassen, ehe er in die Tiefen der Korn-Schlucht gefallen wäre.

Gouri schrie auf. Der Troll gluckste belustigt und rieb sich die Magengegend. Gouri blickte von der Kreatur zu Prip, welcher sich nun an der Aussenseite des Hängebrückengeländers festklammerte, nur wenige Hasensprünge vom sicheren Grasland auf der anderen Brückenseite entfernt... wenn er nur noch kurz durchhielt – doch war Prip sichtlich bedröppelt durch den Aufprall, den er soeben erhalten hatte.

Der Troll gluckste weiter und Gouri fiel viel zu spät auf, wie nahe das Wesen sich nun von ihr befand. Einen langen muskulösen Arm ausstreckend, griff der Troll nach ihr.

Aus dem starken Griff eines Trolls gab es kein Entrinnen. Er hatte die Agren mit einer Hand beinahe komplett umschlossen und hob sie in die Höhe, fast zärtlich, ohne sie zu zerquetschen. Wollte er sich einen Snack gönnen, mitsamt Haut, Haaren und Lederschuhen?

Gouri schrie auf, und Prip stimmte mit ein.

Da liess der Troll sich lachend rückwärts auf den Rücken fallen, sodass die Erde bebte und die Erdgeister rumorten. Wieder im Gras liegend, hielt er die durchgeschüttelte Gouri hoch in die Luft über seiner Körpermitte.

Gouri spürte, dass sich die rauen Finger des Trolls von ihr zu lösen begannen, ehe Prip es sehen konnte – Prip konnte im Dunkeln ohnehin so gut wie nichts sehen. Dann war es auch schon soweit.

Der Troll liess Gouri los.

Sie spürte, wie sie dem vorher so peinigend vorgekommenen Gefängnis entschlüpfte und nach unten glitt. Der freie Fall dauerte nicht lange, doch Gouri kam es wie eine Ewigkeit vor. Der Wind zerrte an ihren verfilzten Haaren und Mooskleidern. Noch im Fall drehte sie sich einmal um die eigene Achse und erreichte den molligen Trollbauch mit dem Kopf zuerst. Die Erfahrung hätte sie unter anderen Umständen als ekelerregend beschrieben, aber besser als ein Besuch in einem Trollmagen war sie allemal.

Gouri prallte mit einem dumpfen Geräusch vom Bauch des Trolls ab und wurde dadurch erneut in die Luft geschleudert, diesmal seitlich. Sich schon auf Knochenbrüche oder schlimmeres einstellend, drehte sie sich, diesmal kontrolliert, sodass sie wenigstens mit ihren Füssen am Boden einschlagen würde.

Sie erreichte den Boden nicht. Stattdessen war des Trolles erdige Hand erneut da, und fasste sie auf. Als er Gouri erneut anhob, erkannte diese an seinem friedlich grinsenden Gesichtsausdruck, dass er ihr nichts Böses wollte. Der tat nix, der wollte nur spielen!

Prip zog sich unter Ächzen über das Brückengeländer und stand auf wackeligen Beinen, mit offenem Mund auf die Szene schauend, die sich vor ihm abspielte.

Der Troll indes liess Gouri erneut los. Ihr Magen drehte sich um und der Trollmagen schlug ihr beim Aufprall die Luft weg, ehe sie erneut in die Höhe geworfen wurde. Diesmal erwartete sie den Griff des fröhlich glucksenden Trolls und bereitete sich innerlich vor, doch streifte der Troll sie unschön am rechten Arm, welcher unangenehm nach hinten gebogen wurden. Es war Zeit, diesem Spielchen ein Ende zu bereiten. Sie holte tief Luft, versuchte, die Schmerzen in ihrem Arm zu ignorieren, und...

„Hey!“, rief Prip von der Brücke aus. Der Troll hob seinen Kopf und starrte ihm grinsend entgegen.

„Hey!“, liess nun auch Gouri verlauten. Der Troll blickte wieder sie an, sichtlich verwirrt. Offenbar verstand er den drohenden Unterton in ihrer zitternden Stimme. Er rappelte sich auf, setzte sich auf seine kurze Unterschenkel und hielt Gouri vor sein Gesicht, welche unter den pochenden Schmerzen in ihrem rechten Arm Mühe hatte, einen wütenden Gesichtsausdruck beizubehalten. Der Troll öffnete sein Maul leicht und Gouri erblickte eine Reihe schiefer gelber Zähne. Wollte er dennoch einen Snack gönnen?

„Hier drüber! Schau hierher!“, hustete Prip hervor. Der lädierte Schildzwerg war von der wackeligen Brücke getreten und etwas näher gehumpelt. Von seiner ersten gewaltsamen Konfrontation mit dem Troll hatte er gelernt, dass er hier mit roher Kraft nichts ausrichten konnte.

Prip hielt die Überreste seines Schildes hoch – hätte genauso gut ein zwergengrosses Stück Holzkohle sein können – und sprach langsam, als würde er sich einem Kind nähern: „Tausch! Tauschen! Dieser Schild hier gehörte einst meinem Vater Aigar, und er war ein mutiger und ehrenvoller Mann, so wie ich voller Ehre bin.“

Der Troll lächelte dümmlich und sein Blick schweifte schon wieder zu Gouri hinüber.

„Schau her!“, warf Prip jetzt ein, „Dieser Schild hat die Feenwelt besucht und ist wieder zurückgekehrt! Feen! Kennst du Feen?“

„Feeeeeeeen“, wiederholte der Troll das Wort langsam. Er konnte ja doch sprechen!

„Feen, genau. Diese kleinen fliegenden Wesen, du kennst sie, oder? Sie sind ausgesprochen lustig, nicht? Dieser Schild ist genauso lustig. Und ich will ihn mit dir tauschen.“

Prip tauschte einen nervösen Blick mit Gouri. Wie viele von seinen Worten verstand der Troll?

„Tau – schen“, brummelte der Troll jetzt.

„Genau, tauschen“, versuchte Prip erneut sein Glück. Er näherte sich langsam dem Troll, während er weiterhin beruhigend auf ihn einredete, „Das heisst: Ich gebe dir was, und du gibst mir was. So haben wir beide etwas davon. Ich schenke dir, du schenkst mir.“

„Dieser Troll geben dir. Du geben etwas an dieser Troll“, brabbelte der Troll nach. Gouri zweifelte inzwischen daran, dass er Prip verstand. Doch der Griff des Trolls hatte sich etwas gelockert, weshalb ihr Arm weniger verrenkt wurde. Das war gut.

Prip pries weiterhin heiser seine kümmerlichen Schildreste an: „Schau nur auf diesen wunderschönen Schild. Die Feenkönigin höchstpersönlich hat ihn einst gesalbt, weil mein Vater die Feenwelt von einigen Dunklen Schemen gereinigt hat. Er...“ Prips Stimme versagte. Gouri sprang ein: „Man sagt diesem Schild nach, er könne bis zum heutigen Tag Übergänge zur Welt der Feen öffnen. Willst du einen Schild besitzen, der Tore zu einer anderen Welt öffnen kann?“

„Feeeeen“, sprach der Troll, und seine Mundwinkel zogen sich noch weiter nach oben, als Gouri je für möglich gehalten hätte.

Prip hatte seine Stimme wiedergefunden und lenkte die Aufmerksamkeit des Trolls wieder zu ihm: „Hier, ich lege den Schild auf diesem Grasstück ab. Aber damit du ihn haben kannst, musst du deine andere Hand öffnen. Dann kannst du zu den Feen reisen! Feen!“

Nun strahlte das Gesicht des Trolles, und er griff mit seiner freien Hand nach den Schildüberresten, die Prip auf dem Boden platziert hatte. Prip reagierte blitzschnell: „NEIN!“

Der Troll zuckte zurück und blickte ihn aus traurigen Augen an. Prip hatte das Tauschgeschäft schon lange aufgegeben, und sprach weiter: „Nein, warte! Zwei Hände! Zwei Hände an den Schild! Sonst keine Feen!“

„Keine Feen?“, sprach der Troll enttäuscht nach. Prip nickte und zeigte ihm beide offenen Hände. Und wie durch ein Wunder tat der Troll es ihm nach und öffnete seine Hand. Gouri purzelte von der Hand hinunter und schlug mit dem Rücken auf dem Pflaster der alten Zwergenstrasse auf. Es schmerzte schrecklich, doch jetzt war keine Zeit dafür, sie musste sich bewegen, schnell, sie war frei! Auf jetzt, auf zur rettenden Festung!

Gouri stellte sich auf ein Knie, und versuchte, sich zu erheben. Die Welt schwankte, doch sie blieb stehen – soll noch einer behaupten, Agren hätten ein schlechtes Gleichgewicht! Sie tastete ihre Gliedmassen und ihr Gesicht ab. Alles brannte unangenehm, und sie schmeckte ein wenig Blut, aber schlimmer als einige Schürfungen würde es nicht sein. Sie humpelte von der riesigen Gestalt des Trolles weg. Ein Blick zurück verriet ihr, dass der Troll die Schildüberreste freudig in die Hände genommen hatte und auf einem Finger zu balancieren versuchte.

Prip nahm nun auch seine Beine in die Hand und huschte hurtig humpelnd am grossen Troll vorbei. Dieser blickte ihm enttäuscht nach und brummelte: „Kein Feen... dieser Schild nicht sein toll! Schild von Troll sein besser. Troll haben Schild mit Händeschütteln drauf! Dieser Schild dort viel mehr toll!“

Prip rammte seine Fersen in den Steinboden der alten Zwergenstrasse und drehte sich rasant um.

„Schild mit Händeschütteln drauf?“, fragte er argwöhnisch nach. Gouri hätte ihm zugezischt, er solle doch einfach still sein, doch sie war noch zu erschöpft dafür, presste ihre Hände in ihre Seiten und verschnaufte tief. Noch schien ja keine imminente Gefahr vom gutmütigen Troll auszugehen.

„Schild mit Händeschütteln drauf!“, nickte der Troll. „Dieser Schild dort sein viel mehr gut!“

Prip verzog sein Gesicht zu einer schmerzverzerrten Grimasse: „Ein rotbrauner Schild?“, fragte er weiter, seine Stimme zitternd.

„Dieser Schild dort haben Farbe von Erdtroll, jawohl.“, nickte der Troll glücklich.

Prip fluchte in seine Bartstummel hinein und drehte sich wieder vom Troll ab, jetzt sogar ein wenig schneller rennend als zuvor. So jagte er an Gouri vorbei, welche sich wieder zu ihm umdrehte und ihm verwirrt hinterhersah.

„Freunde, wartet! Dieser Troll hier können erzählen mehr von dieser Schild dort! Dieser Zwerg dort rennen in Bauch von dieser Troll nochmal! Bitteeee!“

Sein Klagen verklang, als auch Gouri sich in Bewegung setzte und den Troll hinter sich liess. Keine weiteren Geräusche kamen vom Troll. Offenbar war er nicht in Stimmung, ihnen nachzujagen. Zumindest hoffte Gouri dies.

Prip rannte weiter vorne immer noch die Pflastersteine der alten Zwergenstrasse entlang, wurde allerdings etwas langsamer und liess Gouri zu sich aufholen. Gouri schmerze es inzwischen nicht nur an den Prellungen und Schürfungen, sondern auch ein Seitenstechen machte sich bemerkbar. Prip war eindeutig zäher als sie. Warum war er nur so in Eile?

„Warum bist du nur so in Eile?“, presste Gouri zwischen zwei Atemzügen hervor, „Wir könnten hier schnell anhalten, an den Büschen dort drüben wächst etwas Wolfskraut, das könnte ich zerreiben, und die Drei Wasser...“

„Die Drei Wasser sind gefallen, Gouri!“, sprach Prip. „Der Troll beschrieb einen rotbraunen Schild mit zwei verschränkten Händen. Das ist der Bruderschild, einer der vier legendären mächtigen Schilde, die der Meisterschmied Kreatok geschaffen hatte, ehe er seinen Verstand verlor.“

Gouri folgte ihm langsam: „Und wenn dieser Troll sagt, dass er diesen... Bruderschild besitzt...“

Prip vollendete ihren Satz: „...dann ist der Bruderschild sicherlich nicht mehr in den Drei Wassern, wo er sein sollte. Die Drei Wasser sind die Feste, wo die vier mächtigen Schilde Kreatoks aufbewahrt werden. Wenn der Bruderschild daraus errungen werden konnte...“ Er sprach nicht weiter.

„Prip. Wir wissen doch noch nicht mal, ob der Troll die Wahrheit sagt. Vielleicht hat er die Schilde auch nur von weitem gesehen.“

„Die Schilde werden nur in absoluten Notfällen hervorgebracht! Schon allein, dass dieser Troll darüber Bescheid weiss, heisst, dass die Feste angegriffen wurde. Nach dem Fall der Burg Karulzar und des Alten Brons sind die Drei Wasser der letzte oberirdische Zwergenstützpunkt, der uns noch bleibt. Wenn dieser auch noch an die Drachen fällt, werden wir uns endgültig in den Untergrund zurückziehen müssen... doch solange die Drachen da sind, werden wir die Besatzung der Feste nie abziehen können, dann stehen wir da ja wie auf dem Präsentierteller... und im Unterirdischen Kampf sind wir unterlegen, wir brauchen Ballisten, um eine Chance gegen die Drachen zu haben... wenn wir uns in den Untergrund zurückziehen, können die Drachen und die Riesen aus dem Süden sich weiter bekriegen, aber das Volk der Zwerge würde dabei so gut wie ausgelöscht... die Schlacht um die Drei Wasser wird das Schicksal der Schildzwerge bestimmen... ich muss zu den Drei Wassern, ich muss es wissen...“

Prip war zusammengesackt und murmelte ein wenig weiter: „Am Baum der Lieder in den nördlichen Trolllanden sagte man mir, die Riesen aus dem Süden wären kurz davor, die Magie der Toten freizusetzen. Die Toten zurückzuholen! Ich wollte den Fürst der Drei Wasser davor warnen, aber nun...“

Gouri betrachtete den verzweifelten Schildzwerg, nicht wissend, was nun zu tun sein. Auf der Suche nach einer Ablenkung warf sie ein: „Stimmte dies mit dem Schild deines Vaters und der Feenwelt?“

Prip guckte auf und sah sie entrüstet an, sein Blick sprach: „Wie könntest du annehmen, dass ich nicht die Wahrheit sagen würde!?“, während sein Mund sprach: „Es... es ist nur eine Legende, aber genug, um einen Troll abzulenken. Vielleicht erzähle ich dir eines Tages davon.“

Er fasste sich etwas und richtete sich wieder auf. „Ich muss es wissen! Auf zu den Drei Wassern!“

Gouri nickte zaghaft, rupfte ein wenig Wolfskraut aus einem Busch am Wegesrand, steckte es in die Tasche ihres Wams und setzte Prip nach. Gemeinsam huschten sie noch ein kurzes Stück der Zwergenstrasse entlang, welche nun im Mondeslicht deutlich zu erkennen war, da sie den Nebel verlassen hatten. Dann bog Prip nach rechts ab, eine Böschung hoch, und Gouri folgte ihm. Langsam krochen sie bis auf die Spitze des Hügels und blickten darüber.

Hitze schlug in ihre beiden Gesichter

Die Drei Wasser! Auch im nebeligen Dunkel der Nacht konnten Gouris Agrenaugen erkennen, dass es sich dabei um eine wunderschöne Zwergenfeste handelte: Hohe Kuppeltürme, Zwergenstatuen links und rechts, eine Quelle, die in einem riesigen Brunnen sprudelte. Doch das aus dem Quellenbrunnen weit in die Höhe schiessende Wasser konnte nicht verbergen, dass die Festung lichterloh brannte. Flammen schossen meterweit in die Höhe, Zwerge wuselten um den Stein herum und versuchten, das Zusammenfallen des Gebäudes zu verhindern.

Prip riss seine Augen auf: „Wie konnte dies geschehen?! Unser Gestein brennt nicht! Welche diabolischen Mächte – “

Dann fiel sein Blick in die Luft, und er sah die vielen Drachen, die die Feste weit in der Höhe’ umkreisten. Ein einzelner Drache lag weiter weg am Boden, von mehreren dicken Lanzen durchbohrt, doch die anderen hatten sich in die Lüfte geschwungen und blubberten nun fröhlich Lava auf die Drei Wasser hinunter.

Die Drei Wasser brannten lichterloh.

Die letzte Zwergenfeste war gefallen.

Prip schluchzte auf.

„Das... das ist nicht gut.“, murmelte Gouri, den Blick nicht loslösen könnend vom Inferno, welches über den Drei Wassern tobte. Die Flammen tobten, und die hunderte von Metern über den höchsten Zinnen kreisenden Drachen wirkten nicht, als hätten sie vor, den Angriff abzubrechen, ehe von den Drei Wassern mehr als ein verkohltes Gerippe übrig war.

Der gefallene Drache lag verrenkt neben einem der drei hohen Türme, Speere aus seinem Bauch ragend. Offenbar war es den Zwergen möglich gewesen, eine der riesigen Echsen vom Himmel zu holen. Warum versuchten sie es nicht erneut? Sicherlich konnten die Schildzwerge Ballisten fabrizieren, die höher reichten, als die Drachen fliegen konnten?!

Doch ehe Gouri Prip nach diesem Sachverhalt ausfragen konnte, hatte Prip sich erhoben und war den Hang hinuntergestolpert, direkt auf die sterbende Zwergenfeste zu.

Gouri fluchte „Spornendreck!“ und setzte ihm nach.

Glücklicherweise waren die Drachen ziemlich akkurat darin, die Drei Wasser mit Lava und Feuer zu treffen, sodass man ausserhalb der Feste in relativer Sicherheit war. Viele Schildzwerge standen etwas abseits in Gruppen zusammen. Es war nicht zu erkennen, ob sie sich berieten oder einfach nur litten unter dem Anblick ihrer Heimat, die vor ihren Augen davonschmolz.

Prip eilte zu einem solchen Trio. Sowohl auf seinem als auch auf ihren Gesichter erschien ein gegenseitiges Erkennen:

„Prip! Was ist mit deiner Rüstung geschehen?“

„Vergiss deine Rüstung, was ist mit deinem prächtigen Bart geschehen?!“

„Und wer ist denn deine Begleiterin?“

Prip verschnaufte und antwortete: „Das ist Gouri, die mutige Agren. Sie hat mein Leben vor einem dieser Mistviecher gerettet.“

Dann setzte er nach: „Kümmert euch doch nicht um mich! Was in allen Feuern der Tiefminen ist denn hier geschehen?!“

Einer der anderen Schildzwerge, ein grossgewachsener Bursche mit einer Glatze und einem mächtigen roten Schnurrbart, antwortete mürrisch: „Nach was sieht es denn aus? Den Drachen war Karulzar und der alte Bron noch nicht genug! Wir müssen unser Volk nach Cavern evakuieren, aber nicht mal dorthin können wir unbeschadet, da der nächste freie Eingang auf der anderen Seite der Korn-Schlucht liegt!“

„So ein Zufall, von da kommen wir gerade“ liess Gouri verlauten. Der Zwerg mit dem prächtigen roten Schnurrbart blickte sie finster an und fuhr fort:

„Fürst Bailor hat es erwischt. Hornheim aus dem Roteisenstein gibt jetzt den Ton an – wenn es nach diesem Tag überhaupt noch eine Zwergenarmee gibt, über die man herrschen kann. Wir sitzen hier auf dem Präsentierteller!“

Prip fragte entsetzt: „Abes was war mit den Schilden?! Was ist mit den Ballisten?“

Der Zwerg mit dem Schnurrbart murrte: „Du hast einiges verpasst. Die Drachen haben uns belagert, die Ballisten hervorgelockt, uns unsere Speere verschwenden lassen. Mit Verlaub, unser guter Kommandant war ein Trottel. Und die Nachschübe aus Cavern haben auf sich warten lassen. Wahrscheinlich haben die elenden Riesenreptilien Feuerstösse ausgelöst und uns vom Rest des Höhlennetzwerks abgeschnitten.“

Der Zwerg holte tief Luft und fuhr fort: „Sie haben die Wachtürme und Waffenkammern zuerst abgesengt. Das war schon einige Tage her. Als hätten diese steinernen Trottel es gerochen, ist eine Horde Trolle gekommen und hat sich den Bruderschild aus den Überresten der Waffenkammer gekrallt, ehe wir einen Gegenschlag ausführen konnten. Gut, wir waren auch anderweitig beschäftigt.“

Gouri wusste besser, als den Troll von vorhin mit dem Bruderschild zu erwähnen. Der Zwerg war schon enerviert genug.

Sein Schnurrbart erzitterte, als der Zwerg fortfuhr: „Den Sternenschild und den Dunkelschild konnten wir den Trollen entreissen, doch fiel der Sternenschild kurz darauf an die Drachen, die im Sturzflug Zwerg um Zwerg von unseren Zinnen picken. Es sind mehr gefallen als damals, als der Urtroll Tiefenfall vernichtete! Das ist das Ende der Schildzwerge!“

Nun sprach der Zwerg nicht mehr. Prip und Gouri warteten immer noch auf den Aufenthaltsort des Dunkelschilds.

Ein anderer Zwerg meldete sich zu Wort: „Tja, der Dunkelschild... der liegt dort drüben unter diesem Biest!“ und er zeigte zum gefallenen Drachen, welcher nahe der Mauer der Drei Wasser am Boden lag.

Die Zwerge fluchten gemeinsam. Unter diesem Vieh würde man den Schild nur sehr mühsam und unter grossem Zeitaufwand bergen können. Zeit, die sie schlichtweg nicht hatten.

Gouri, die keine Ahnung vom Armeegeschäft hatte, aber dutzende Legenden über die mächtigen Schilde kannte, konnte verstehen, was für ein riesiger Verlust das Fehlen der drei Schilde für die Schildzwerge war. Gute Güte, vielleicht würde das wirklich das Ende der Schildzwerge sein.

„Immerhin, wir sind dabei, einige Wurfmaschinen zu bauen. Wenn diese Sklavenschinder aus dem Süden Zwerge dazu anleiten können, ganze Zeras zu bauen, so können wir freien Schildzwerge das erst recht! Wurfmaschinen mögen weniger elegant als Ballisten sein, aber sie werden ihre Arbeit tun. Und die Drachen konzentrieren sich bloss auf die Drei Wasser, die werden das Ganze ganz sicher nicht kommen sehen.“

Da atmete Prip wieder ein wenig auf.\bigskip



Die Sonne erhob sich hinter den Bergen, und Gouri verspürte Hunger in ihrem Magen. Sie nahm nicht an, dass einer der an- und niedergeschlagenen Zwerge sich zu einem Morgenmahl begeben würde, und so kletterte sie einen Hügel hinunter und schnappte sie sich einen angrenzenden Busch und begann, ihn knabbernd zu verzehren. Dann blickte sie zu den Zwergen auf, die sie mit grossen Augen ansahen.

„Oh Herrin des Steins, das war doch hoffentlich kein heiliger Busch oder so was?“, rief Gouri ihnen entgegen. Die Zwerge starrten immer noch, und Gouri war sich immer noch sicher, dass es mit dem Verzehren ihres Asts zu tun hatte. Doch dann lachten sie auf und drehten sich weg, weiter an einem verzweifelten Rettungsplan für die überrannte Zwergenfeste schmiedend.

In diesem Moment vernahm Gouri ein Knacken aus dem Gebüsch hinter ihr.

Sie drehte sich alamiert um.

Eine verhutzelte Gestalt krabbelte kaum zehn Meter von ihr entfernt durch das Gestrüpp. Etwas weiter entfernt erkannte Gouri etwas, das an eine Ruine erinnerte – dort musste einst ein Aussenturm der Drei Wasser gestanden haben, welcher jetzt natürlich verbrannt und eingestürzt war.

Wahrscheinlich hatte die Person vor ihr sich im Turm befunden, als dieser gestürzt worden war, und nun hatte sie sich endlich in die Freiheit gekämpft.

Die Gestalt war ein Zwerg, zweifelsohne, mit einem langen grauen Bart. Er wäre bestimmt elegant gekleidet gewesen, wäre sein langer Mantel nicht an so vielen Stellen gerissen, verschmutzt und angebrannt.

Etwas an der Erscheinung dieses Zwergs liess Gouris Alarmglocken klingeln. Äusserlich konnte Gouri keine Verletzungen an ihm feststellen, doch ging er langsam und ungelenk, als hätte er seine Füsse für eine lange Zeit nicht mehr gebraucht. Seine Lippen bewegten sich rasch, als murmle er vor sich hin. Fahrig rieb er sich seine Arme.

Da fiel Gouri auf, dass seine Hände und Füsse in den Überresten goldener Ketten eingeschlossen waren. Handelte es sich bei diesem Zwerg um einen Gefangenen der Schildzwerge? Sie hatte gehört, dass die Schildzwerge manchmal spezielle Häuser für diejenigen einrichteten, die gegen allgemein anerkannte Regeln verstossen hatten, und wenn sie einmal drin waren, durften sie sie nicht mehr verlassen.

Gouri stellte fest, dass der entflohene Zwerg in den goldenen Ketten ihr eindeutig zu nahe war für ihren Geschmack. Zwar hatte er noch nicht zu erkennen geben, dass er Gouri bemerkt hatte, doch bewegte er sich in leichtem Zick-Zack auf sie zu, und wenn sie sich nicht irrte, wurden seine Schritte sogar noch schneller. Nein, jetzt war er noch fünf Schritte von ihr entfernt und kämpfte sich buchstäblich durch das Gebüsch auf sie zu. Gouri drehte sich ab und blickte oben an den Hügel, wo Prip und die übrigen Schildzwerge sich berieten. Sie musste so schnell wie möglich zu ihnen.

Die Schildzwerge standen immer noch oberhalb des Hügels, doch sie schienen sich nicht mehr zu beraten, stattdessen blickten sie allesamt auf Gouri, oder genauer gesagt auf einen Punkt ein wenig hinter Gouri. Sie mussten den entflohenen Zwerg ebenfalls gesehen haben.

Prips Gesichtsauszug zeigte ein erschrecktes Erkennen des Zwergs, und er rief aus: „WEG VON IHM, GOURI, WEG!“

Das war alles, was sie zu hören brauchte. Etwas war mit diesem geflohenen Zwerg nicht in Ordnung. Sie musste weg von ihm. Sie musste weg, JETZT!

Und Gouri rannte los, so schnell es ihre müden verkratzten Gliedmassen zuliessen. Nach nur einigen stolpernden Ansätzen wurde sie von hinten gepackt und zu Boden geworfen. Ihr Kopf schmerzte, doch sie konnte nicht ausruhen, nicht jetzt!

Ein überraschend starker Griff packte sie am Kragen ihrer Mooskleidung und zog sie nach oben. Der entflohene Zwerg hielt sie fest vor sich und blickte ihr in die Augen. Seine eigenen Augen zuckten unwillkürlich, und er rief aufgeregt: „Ich habe ihn nicht getötet, weisst du, ich habe ihn nicht getötet, und dass er nicht zurückkommt, nie mehr zurückkommt, das ist seine Schuld, nicht meine!“

Gouri versuchte zurückzuweichen, dem Griff zu entkommen, aber der Zwerg zerrte sie nur noch weiter in die Höhe, sodass ihre Füsse kurzzeitig den Boden verliessen.

„So versteh’ mich doch, ich wollte das nie! Die haben mir gesagt, das könne ich nicht mehr wiedergutmachen, nur dafür büssen! Büssen! Habe ich nicht schon genug gebüsst?! O Nehal, warum hast du mich nur verlassen?“

Dann lockerte sein Griff sich und Gouri plumpste zu Boden. Um sich blickend, erkannte sie Prip und die drei weiteren Schildzwerge, die näher getreten waren und einen Kreis um den entflohenen Zwerg und die am Boden liegende Gouri bildeten. Sie waren dem geflohenen Zwerg zahlenmässig weit überlegen, doch waren alle ihre Gesichter angespannt. Selbst Prip warf nur einen raschen Blick auf Gouri, um zu sehen, ob es ihr vergleichsweise gut ging, und blickte dann blitzschnell zurück zum Zwerg in den Ketten.

Dieser grinste: „Vorsicht ist die Mutter aller Tugend, das hat Nehal immer gesagt. Ich nehme dann an, ihr wisst, wer ich bin? Ich...“, er lachte nervös auf, „Ich bin mir nämlich gerade nicht so sicher, wer ich bin. Ich weiss natürlich, wer ihr denkt, dass ich bin...“

Die Zwerge tauschten unsichere Blicke aus. Der mit dem prächtigen Schnurrbart nickte grimmig.

„Aha!“, rief der Zwerg in den goldenen Ketten aus, „Ihr glaubt, zu wissen, wer ich bin! Dann solltet ihr ja auch wissen, dass ich bereits gebüsst habe, ihr Narren! Und ich gehe sicherlich nicht mehr zurück in dieses Drecksloch! Also... ihr werdet mich ziehen lassen müssen, oder ich sehe mich gezwungen, zu... drastischeren Massnahmen zu greifen.“

Während die Sprechweise des Zwergs kohärenter wurde, wurde auch sein Blick klarer, und er schien sich seiner Situation besser bewusst zu werden. Das war nicht beruhigend. Er drehte sich einmal langsam im Kreis und musterte jeden Zwerg kurz. Bei Prip blieb sein Blick hängen:

„Was ist denn mit dir passiert?“

„Das war natürlich ein Drache. Keiner von denen hier. Und er ist jetzt tot“, erwiderte Prip grimmig. Gouri konnte nicht erkennen, ob er versuchte, Eindruck zu schinden.

Beeindruckt wirkte der entflohene Zwerg nicht, stattdessen liess er die Überreste seiner goldenen Ketten locker um seinen Arm herumwirbeln, und bellte: „Ein Drache, soso. Wie hiess er denn?“

„Sagrak.“

Der entflohene Zwerg grinste: „Oho, das freut mich doch. Weil ich jemanden kenne, den das ganz und gar nicht freuen wird.“

Dann schrie er plötzlich: „HA!“ und griff an seinen linken Unterarm. An diesem Handgelenk war eine kleine, surrende, golden glänzende Apparatur befestigt, mit einer Sprungfeder, welche sich soeben von einem Haken löste – und ein goldener Bolzen schoss aus der Apparatur hervor und traf der schnurrbärtigen Zwerg an der Schulter. Er verzog das Gesicht schmerzverzerrt und fiel zu Boden. Prip und die beiden anderen Schildzwerge – ein breiter Krieger und eine hochgewachsene Kriegerin – brauchten einige Sekunden, um sich zu fassen. Sekunden, die der entflohene Zwerg gut zu wissen nutzte. Er rammte Prip in die Seite, warf ihn zu Boden und gab zünftig Fersengeld, den Hügel hinauf, weg von der Gruppe.

Überrumpelt, aber nicht weiter verletzt, rappelte Prip sich wieder auf und blickte hinüber zum schnurrbärtigen Zwerg, in dessen Schulter immer noch ein goldener Bolzen steckte – der entflohene Zwerg hatte genau in die Spalte der Rüstung getroffen.

Der breite Krieger und die hochgewachsene Kriegerin hatten sich bereits neben ihm niedergelassen:

„Wie schlimm ist es?”

„Meinst du, der Bolzen ist vergiftet?”

„Definitiv, sonst würde unser Belenor das doch mit links wegstecken.”

„Aber wie hätte er an Gift kommen können?”

„Wie er wohl an diese Waffe gekommen ist – ach, die hat er bestimmt wieder selbst entworfen. Wie immer! Dieser verflixte Genius.”

Der verletzte Zwerg mit dem Schnurrbart – Belenor – stöhnte auf und die Aufmerksamkeit fiel wieder auf ihn. Nun gelang auch Prip dazu. Gouri hatte bereits ein wenig Wolfskraut aus ihrem Mooskleid gezogen und fragte die übrigen Zwerge: „Will denn niemand den entflohenen Zwerg verflogen?”

Die übrigen blickten sich unsicher an. Dann ergriff Prip das Wort:

„Vielleicht… vielleicht ist das gerade das, was er will.”

„Wir können nicht wissen, was er vorhat. Wenn er in so kurzer Zeit eine kleine Giftmischung herstellen konnte, hat er inzwischen schon eine geniale Falle erbaut”, half ihm die Zwergin aus.

„Vielleicht ist es das Beste, ihm komplett aus dem Weg zu gehen”, hustete nun auch Belenor.

Gouri blickte den Schildzwergen einem nach dem anderen ins Gesicht, und lachte dann beinahe auf. Sie hatten Angst! Dann fiel ihr auf, dass sie demnach selbst auch Angst haben sollte.

„Wer... wer war dieser entflohene Zwerg denn?” wollte Gouri wissen.

Prip setzte zu einer Antwort an, doch da ertönte ein langgezogener Schrei von der anderen Seite des Hügels, und die Zwerge blickten einander alarmiert an. Der breite Zwerg fasste sich zuerst und hetzte den Hügel hinauf, gefolgt von Prip und der Kriegerin. Belenor blieb liegend zurück und schien inzwischen nicht einmal so unzufrieden mit seiner Position zu sein. Gouri sah, dass es ihm vergleichsweise gut ging, also liess sie ihr Kräuterbündel bei ihm und spurtete selbst den übrigen Schildzwergen hinterher.

An der Hügelkuppe angekommen, erblickte sie ein grauenvolles Bild. Ein weiterer Drache war vom Himmel gefallen und lag nun am Ende einer tiefen Schleifspur in der Erde, mitten in der Ebene vor der Festung. Er rührte sich nicht mehr. Er schaute auch nicht mehr wie ein lebendiger Drache drein – vielmehr war er steinern in einer fliegenden Position eingefroren.

Keine Zwergenballiste hatte den Drachen erledigt, und keine Lanze steckte in ihm. Gouri blickte wild um sich, konnte aber nichts erkennen, was das Wesen niedergestreckt haben könnte. War er einfach so im Flug versteinert worden?

Schliesslich ergriff Gouri das Wort: „Wer hat denn den da erledigt?“

Prip lachte auf: „Das war er selbst. Weißt du, die Drachen sind Seelen der Erde und des Feuers, zusammengeschmiedet während ihrer Entwicklung im Ei. Doch sie verhalten sich gerne wie Seelen der Lüfte und schwingen sich in ungeahnte Höhen. Werden sie zu übermütig und verbringen zu wenig Zeit in Kontakt mit der guten alten Mutter Erde, so holt diese sie sich auf ihre eigene Art zurück.“

Die hochgewachsene Zwergenkriegerin schnaubte auf: „Als ob die Mutter der Erde so aktiv in das Weltgeschehen eingreifen würde. Das ist doch eindeutig ein Fluch der Windgeister, der allzu aufmüpfigen Drachen zeigen soll, wo sie hingehören!“

Der breite Zwergenkrieger gluckste und meinte: „Mein Vater hat mir jeweils erzählt, dass die Trolle vom Stein kommen, und jedes Mal, wenn der Urtroll einen neuen erschafft, so verlangt die Balance der Natur, dass ein anderes Wesen wieder zu Stein wird, und das ist dann halt eben ein Drache wie dieser gefallene Kerl hier.“

Prip seufzte und argumentierte nicht weiter, sondern setzte einen poetischen Schlussstrich: „Vom Steine sind sie gekommen, und zu Steine sollen sie werden.“

Die Zwergenkriegerin nickte: „Mehr Gestein für unsere zukünftigen Katapulte. Das mag ich.“

Weiter vorne erblickte Gouri den entflohenen Zwerg, wie er mutig über das Schlachtfeld sprintete. Keiner beachtete ihn, alle Augen waren auf den versteinerten Überresten des gefallenen Drachen gerichtet. Sie machte die anderen Zwerge auf den fliehenden Zwerg aufmerksam.

Prip fluchte: „Schaut euch mal seine Trajektorie an. Der will zum Drachenkadaver.“

„Dort liegt der Dunkelschild!“, hauchte Gouri.

Der breite Zwerg nickte zustimmend: „Unter mehreren Tonnen Drachenfleisch sollte der Dunkelschild eigentlich auch vor ihm sicher sein – aber wenn irgendjemand darankommen kann, dann ist es er.“

„Dann müssen wir ihn aufhalten!“

Die Zwerge standen allen Anschein nach einer Konfrontation mit dem entflohenen Zwerg immer noch abweisend gegenüber. Überraschenderweise sahen sie aber nicht bedrückt aus. Die hochgewachsene Kriegerin grinste sogar schwach und meinte dann zu Gouri: „Er kommt nicht weit. Siehst du, wie die übrigen Zwerge aus dem Weg laufen? Die wissen, was da kommen wird.“

Gouri blickte, und tatsächlich! Die meisten Schildzwerge, die sich noch in der Nähe des Gemäuers befunden hatten, huschten von der Einschlagstelle des gefallenen versteinerten Drachen weg, als ginge es um ihr Leben – was darauf hindeutete, dass es tatsächlich um ihr Leben gehen könnte. Einzig der entflohene Zwerg kümmerte sich nicht darum und rannte schnurstracks an den versteinerten Überresten der Echse vorbei, auf den Leichnam des anderen Drachen hinzu.

Prip fuhr fort: „Das hat damit zu tun, dass die Drachen ihre versteinerten Kameraden nicht zurücklassen. Schliesslich können sie sie noch retten, wenn sie sie rasch genug in ihr stinkendes Nest zurückbringen. Warte darauf, was jetzt kommt.“

Erwartungsvoll blickte Prip nach oben. Eine dicke Wolkendecke schwebte inzwischen über den Drei Wassern, aber die immer noch von oben herabfallende Lava zeigte deutlich an, dass die Drachen weiterhin über dem Gemäuer schwebten und Feuer herabregnen liessen. Kam es ihr nur so vor, oder war die Lava weniger geworden?

Da wurde die Wolkenwand zerfetzt, und Gouri fühlte sich ungut an Sagraks Angriff auf Prip erinnert, als ein riesiger schwarzer Drache aus dem Himmel stob und ohne abzubremsen neben dem gefallenen Steindrachen auf die Erde prallte, sodass Schlamm und Gestein hunderte von Zwergenhöhen gen oben und zur Seite stieben. Dann erhob er sich langsam aus dem soeben verursachten Krater.

Sagrak war der erste Drache gewesen, den Gouri gesehen hatte, und schon er hatte ihr eine Todesangst eingejagt. Dieser hier war um ein Vielfaches schlimmer. Er hätte Sagrak um mehr als das fünffache überragt, seine Krallen alleine waren schon fast so lange wie ein mächtiger Zwerg. Sein Schwanz zuckte umher, als er seinen schlangenhaften Kopf um sich wandern liess, mit seinen riesigen rot glühenden Augen auf der Suche nach etwas, das ihm gefährlich werden konnte. Doch da war nichts. Dieser schwarze Drache war der sichere Tod für alle, die es wagten, sich ihm entgegenzustellen.


\newpage
\section{Kreatok gegen Tarok}

Der soeben gelandete schwarze Drache musterte seine Umgebung aufmerksam. Diejenigen Schildzwerge, die durch den Einschlag des Reptils in die Erde noch nicht von den Füssen geholt worden waren, duckten sich hinter ihre Schilde und begannen, alte Zwergenlieder zu rezitieren. Offenbar bereiteten sie sich auf den Mentalen Angriff des soeben gelandeten Riesendrachen vor. Langsam wurden die verschiedenen Stimmen synchron, und mit anschwellender Lautstärke sangen die Zwerge eine wortlose Melodie, in Bruderschaft tapfer vor einem übermächtigen Feind stehend und dessen hinterhältige geistige Attacke erwartend. Sie sangen die uralte Melodie des Steinsang.

Der entflohene Zwerg in den goldenen Ketten drehte nur kurz seinen Kopf, um die Anwesenheit des riesigen Gegners anzuerkennen, und rannte dann schnur stracks weiter auf den Drachenkadaver zu, wo der Dunkelschild lag.

Der riesige Drache verzog seinen Mund zu etwas, das man mit etwas Fantasie ein Grinsen hätte nennen können, und entblösste dabei gelbe Zähne, grösser als Langschwerter. Seine Stimme war tief und höhnisch, als sie in Gouris Kopf erscholl:

„WICHTE! WER DENKT IHR, DASS ICH BIN? ICH HABE ES NICHT NÖTIG, ZU SOLCH EHRENLOSEN MITTELN WIE EINEM GEISTIGEN ANGRIFF ZU GREIFEN. ICH BIN EIN DRACHE VON EHRE. DENJENIGEN UNTER EUCH, DIE IHREM GERECHTEN ENDE STANDHAFT GEGENÜBERTRETEN WOLLEN, GEWÄHRE ICH DEN EHRENHAFTEN TOD DURCH DAS LEUTERNDE FEUER. LASSET EURE SCHILDE FALLEN.“

In Gouris Kopf dröhnte es, und ein rascher Blick auf ihre Begleiter verriet ihr, dass es ihnen ebenfalls so ging. Kein einziger Schildzwerg liess seine Deckung fallen, vielmehr wurde der gemeinsame Gesang nur noch lauter. Der riesige schwarze Drache fuhr unberührt fort:

„WUSSTE ICH’S DOCH! NICHTS ALS WIRBELLOSER SPRONENDRECK SEID IHR, WÜRMER UNTER MEINEN TATZEN. IHR FÜRCHTET EUCH ZU RECHT, DENN ICH BIN TAROK, DER RÄCHER, UND ICH HABE GESCHWOREN, EUCH ALLE FÜR EURE VERBRECHEN BÜSSEN ZU LASSEN!“

Sein Hals richtete sich auf, und Gouri, die die Anzeichen inzwischen erkennen konnte, duckte sich hinter den Hügel. Tarok spie einen gewaltigen Feuerstrahl und tauchte seine Umgebung in ein Flammenmeer. Gouri fühlte die Hitze an ihr vorbeiwallen, und sie duckte sich noch tiefer in die kühle Erde. Geschrei ertönte irgendwo vor ihr, grauenvolles Geschrei. Ein abrupter Schmerz in ihrem Rücken verriet ihr, dass ihre Mooskleidung zu glühen begann. Rasch rollte sie über die Erde, um das Feuer zu löschen, und der Schmerz verklang etwas. Sie beschloss, noch etwas weiter den Hügel hinunterzurollen, weg von dem Flammeninferno.

Da prallte sie mit jemandem zusammen. Sie blickte auf. Hitze schlug ihr in die Augen, sodass sie diese gleich wieder schliessen musste, doch erkannte sie in diesem Lidschlag Prip und die übrigen Schildzwerge, die ebenfalls unten am Hügel das Flammenmeer ausharrten.

„Tarok, der Rächer?“, hauchte Gouri entsetzt. Sie hatte grausame Geschichten über diesen jähzornigen Drachen gehört.

Die umliegenden Zwerge nickten verzweifelt. Belenor war es schliesslich, der mit zitterndem Schnurrbart das Wort ergriff: „Es ist vorbei. Die Furcht vor ausgeklügelten Fallen im Erdboden war es, der die Drachen in der Luft behielt. Jetzt, wo Tarok hier gelandet ist, und unsere Machtlosigkeit demonstriert, wird ihnen klar werden, dass wir ihnen komplett unterlegen sind. Alle Drachen werden hier landen, und die Drei Wasser vernichten. Sie werden die Feste aus dem Erdboden reissen und in die Tiefen von Cavern eindringen. Wir haben verloren.“

Die übrigen Zwerge nickten traurig. Einzig Prip schien zwar fassungslos, den Mut aber noch nicht verloren zu haben: „Jetzt kommt schon! Tarok spricht hier zwar grosse Worte, aber das dient auch nur unserer Abschreckung. Er will sich bloss seinen versteinerten Freund schnappen. Noch ist nicht alles vorbei!“

Er rüttelte den breiten Krieger an den Schultern: „Wir sind viele! Die Drei Wasser werden fallen, doch wir Zwerge werden uns verstreuen, und wie der Phönix aus der Asche werden wir ein neues Reich errichten. Es gibt immer noch die Eier!“

Gouri blickte wieder auf: „Die Eier?“

Prip nickte mit einem schelmischen Grinsen: „Diese elenden Biester müssen ihre Eier tief unter der Erde vergraben, damit sie schlüpfen. Darum bauen sie grosse Bauten und Höhlen. Aber wir Zwerge sind die Herren des Gesteins, und selbst wenn diese zu gross gewachsenen Eidechsen unsere Gänge mit ihrem heissem Atem versengen, bis sie spucken müssen – solange noch ein einziger wahrer Schildzwerg übrig ist, so werden wir weiterhin die unterirdischen Nester der Drachen ausfindig machen und vernichten.“

Die hochgewachsene Kriegerin warf nun ein: „Wir können mit Stolz verkünden, dass seit dem Ausbruch des Kriegs kein einziges Drachenküken mehr das Licht der Welt erblickt hat. Die hätten uns halt nicht verraten und angreifen dürfen!“

Jetzt war Gouri verwirrt: „War es nicht ein Zwerg, welcher den Krieg ausgelöst hatte?“

Die hochgewachsene Kriegerin schluckte tief und meinte: „Theoretisch ja, es war der Fallenmeister und Meisterschmied Kreatok, der ein Massaker unter den Drachen anrichtete. Aber dem hat es ja auch den Verstand zugenebelt, und wir haben ihn vor Gericht gestellt und zünftig lange eingekerkert. Kreatok ist keiner von uns mehr. Die Drachen jedoch, diese elenden Biester, die haben das als Verrat der Schildzwerge an ihrem Volke interpretiert und sich gegen uns gewendet. Der Rest hier ist reine Selbstverteidigung. Die Schildzwerge und Verrat, was für eine Vorstellung.“ Sie spuckte angewidert auf den Boden.

Gouri blickte auf, überrascht bemerkend, dass sie die Hitze nicht mehr spürte. Das Flammenmeer war versiegt. Die Erde zitterte nicht mehr. Sie guckte nach oben, doch die Hügelkuppe verdeckte die Sicht auf das Schlachtfeld. Hatte Tarok bereits den versteinerten Drachen geschnappt und in Sicherheit gebracht? War etwas vorgefallen?

Da ertönte ein metallischer Klang, eine Mischung aus dem dumpfen Schallen eines Gongs und dem zischenden Erwachen einer Feuerlunte. Gouri schaute verwirrt um sich. Sie konnte nicht erkennen, von wo der Klang stammte.

Auf den Gesichtern der Zwerge tauchten jedoch breite Grinsen auf. Sie erkannten diesen Klang. Prip stemmte sich sogar ohne zu zögern hoch und rannte zur Spitze des Hügels. Gouri wollte ihn anschreien, er solle unten bleiben, doch dann zogen auch die übrigen Zwerge an ihr vorbei, ungeachtet der Gefahr, die immer noch von Tarok ausgehen könnte. Gouri haderte für einen Augenblick, und rappelte sich dann ebenfalls auf, den Hügel hinaufhastend. Oben stehend hatte sie einen guten Blick über das Schlachtfeld. Und was für ein Anblick das war!

Tarok, der riesige Drache, stand immer noch in dem von ihm verursachten Krater neben dem versteinerten Drachen. Rund um ihn war die Erde schwarz und verdorrt, die Gräser verschmort und die Steine gesplittert. Den meisten Zwergen, die sich hinter ihre Schilder geduckt hatten, war es nicht besser ergangen, doch hatte isolierende Kleidung, wie Prip sie trug, vielen das Leben gerettet, und nun sassen sie ächzend in den angeschmolzenen Überresten ihrer Rüstung fest.

Die Drei Wasser standen auch immer noch, und immer noch regnete es ein wenig Lava von der Drachenhorde oberhalb der Wolken.

Doch der Drachenkadaver war nicht mehr. Wo er gelegen hatte, wogte jetzt ein Meer aus schwarzen und silbernen Flammen, die den Drachenkörper bereits verzehrt hatten und an der umgebenden Erde zu nagen begannen. Dieses Flammenmeer war fast so gross wie Tarok selbst, und der Drache fixierte das fremdartige magische Gebilde mit verengten roten Augen, sein Schwanz immer stärker zuckend. War er wütend? War das Furcht? Gouri konnte es nicht sagen.

Und da erkannte sie endlich die Quelle des seltsamen metallenen Geräuschs. Vor dem silbernen Flammenkonstrukt stand ein Zwerg. Nicht irgendein Zwerg, nein, es war derselbe entflohene Zwerg mit den Überresten goldener Ketten an seinen Händen und Füssen, welcher vorhin ihre Gruppe angegriffen und Belenor verletzt hatte. Jetzt stand er mutig da, und stellte sich dem Drachen entgegen. In seiner Hand hielt einen spitzen schwarzen Schild mit wenigen Ornamenten am Rande, auf den er immer wieder mit seiner goldenen Kette schlug, was das metallische Klingen auslöste. Das silberne Feuer wogte und waberte im Einklang mit seinen Schlägen. Kein Zweifel: Das war der legendäre Dunkelschild, der den Unterirdischen Krieg ausgelöst hatte. Der entflohene Zwerg war der Träger des Dunkelschilds. Der Erschaffer des Dunkelschilds. Der entflohene Zwerg war der legendäre Meisterschmied Kreatok. Und er war soeben mit seiner gefährlichsten Kreation wiedervereint worden.

Wie der verwirrte Zwerg, welcher vor Kurzem auf Gouri zugestolpert war und wirr vor sich hin geredet hatte, wirkte Kreatok nun wahrlich nicht mehr. Sein Blick war wieder fest nach oben auf Tarok gerichtet, ein hämischen Grinsen unter dem mächtigen Rauschebart erkennbar. Dann hörte er auf, weiteren Lärm mit seinem Dunkelschild zu machen, und öffnete seinen Mund zu einem Triumphgegröle:

„Meine Brüder und Schwestern, wie sehr habe ich eure Gesellschaft vermisst! Ich bin der Kreatok, den ihr viel zu lange falsch behandelt habt – doch fürchtet euch nicht, ich bin nicht nachtragend. Unter meiner Abwesenheit hat das Volk der Schildzwerge gelitten, so wie auch ich gelitten habe, alleine eingesperrt in diesem Turm, weit weg von meinen schönen Kreationen. Doch frohlocket und jauchzet, denn ich bin zurück, und somit wird alles wieder in Ordnung sein!“

Es war eindeutig nicht alles in Ordnung. Gouri schluckte schwer. Sie wollte sich nach den übrigen Zwergen umsehen, insbesondere Prip, um zu sehen, was sie von dieser Entwicklung hielten. Klar, Kreatok war ein Verbrecher, und ein verrückter noch dazu, doch schien er in Verbindung mit den mächtigen Schilden eine wahre Chance gegen die übermächtigen Drachen zu bieten. Oder war das nur Farce? Er wirkte weniger wahnsinnig als vorhin, doch für wie lange würde das so bleiben? All diese Fragen erhoffte Gouri sich von Prip beantwortet zu kriegen.

Doch Prip war nicht mehr hier.

Die übrigen Zwerge waren nicht mehr hier.

Gouri erkannte sie weiter vorne, bereits tapfer aufs Schlachtfeld rennend, Prip ungeachtet seiner zahlreichen Wunden und blauer Flecken. Sie wollten ihren Freunden helfen.

Von den Rändern der Ebene und aus dem Innern der der Drei Wasser erkannte Gouri weitere Zwerge, die hervorströmten, zu den Gefallenen und Verwundeten, teils mit Tragen, teils mit Heilkräutern, teils mit nichts ausser ihren Kleidern und dem Willen, die verwundeten und angekokelten Schildzwerge in Sicherheit zu tragen.

Das war auch Kreatok aufgefallen, und er schrie weiter: „Ja, meine Brüder und Schwestern, schnappt euch die Verwundeten und bringt sie in Sicherheit. Ich bin hier, um mich wieder in den Dienst meines Volkes zu stellen. Wenn all das hier vorbei ist, dann werde ich hinunter in die Hallen von Cavern steigen, den Thron beanspruchen und die Zwerge der mächtigen Schilde in ein glorreiches Zeitalter der Innovation bringen! Ihr habt zu lange im Schatten der Riesenechsen gewartet, und zu viele von uns leiden IN DIESEM AUGENBLICK unter dem Joch der Riesen aus dem Süden! Wir werden sie alle befreien! Doch zunächst gibt es eine letzte Riesenechse, die mich enttäuscht hat, und mit der es fertigzuwerden gilt.“

Damit richtete Kreatok seine starren Augen wieder auf Tarok, welcher seit dem Erscheinen des Schwarzen Feuers keinen Wank getan hatte: „Tarok, du alter Waschlappen! Ich habe gehört, dein kleiner Sagrak ist nicht mehr nach Hause zurückgekehrt. Das muss schon tragisch sein, auch noch deinen letzten Sohn zu verlieren. Echt, wenn du lieber in dein heimeliges Nest gehen und dich dort ausweinen willst, lasse ich dich bereitwillig ziehen, mein Freund!“

Taroks Gesicht zeigte keine Regung, aber eine Welle der Wut und des Schmerzes brandete abrupt von ihm aus. Gouri war überrascht, wie gut Kreatok Tarok zu kennen schien. Und darüber, dass Sagrak, der Drache, welcher gestern Prip überfallen hatte und von einem Troll getötet worden war, offenbar Taroks Sohn gewesen war. Das machte die Situation auf keinen Fall besser.

„Hast du deine Worte verloren? Schon klar, es kommt nicht jeden Tag, dass man...“, setzte Kreatok erneut an. Doch diesmal unterbrach Tarok seinen Spott und sprach langsam und deutlich ein einziges Wort:

„KREATOK.“

Kurzzeitig flackerte Unsicherheit in Keratoks Stimme auf, ehe er zu einer Antwort ansetzte: „Ja, so haben mich meine Eltern benannt. Worauf willst du hin...“

Erneut wurde er von Tarok unterbrochen, welcher tonlos sprach: „ICH HÄTTE NICHT GEGLAUBT, DASS DU NOCH AM LEBEN BIST. DEINE ZUNGE IST SPITZER GEWORDEN.“

Er wartete einen kurzen Augenblick, und setzte dann nach: „NEHAL VERMISST DICH.“

Gouri konnte von ihrer entfernten Position Kreatoks Gesicht nicht erkennen, doch es kam keine Antwort von Seiten des Zwergs. Dafür flammte das Schwarze Feuer hinter ihm weiter auf, sodass es Tarok nun deutlich überragte. Nehal war der Name von Kreatoks besten Freund gewesen, einem Drachen. Kreatok hatte ihn der Legende nach mit ebenjenem Schild umgebracht, den er nun wieder in seiner Hand hielt.

„KREATOK. DEINE TATEN SIND UNVERZEIHLICH, DARUM SCHWÖRE ICH DIR GLEICH JETZT: SOBALD ICH DIE GELEGENHEIT DAZU KRIEGE, WERDE ICH DICH VERNICHTEN. ABER DEIN VOLK MUSS NICHT WEITER UNNÖTIG LEIDEN. ICH BIN EIN DRACHE VON EHRE. WENN DU DICH ERGIBST, SO LASSE ICH DIE ANDEREN ZIEHEN.“

Kreatok sagte nichts und die Flammen hinter ihm wuchsen weiter in die Höhe und Breite, erreichten die Drei Wasser und leckten an der bröckelnden Aussenmauer.

„KREATOK, SEI VERNÜNFTIG. DEINE ZEIT IST UM. NEHAL HAT DIR SCHON LANGE VERGEBEN. ER WARTET AUF DICH AUF DER ANDEREN SEITE.“

„Es war seine Schuld!“, schrie Kreatok jetzt plötzlich auf. „Er soll mir nicht verzeihen, das lag alles an ihm! Ihm hätte das Feuer nichts angehabt, hätte er sich nicht von mir abgewandt! Ich habe ihn nicht... ich wollte nicht!“

„WAS WOLLTEST DU NICHT, KREATOK?“

„Ich wollte keinen Krieg! Es gab nie ein geheimes Bündnis der Zwerge gegen die Drachen. Ich teilte mein Wissen über die Dracheneier nur aus Notwendigkeit... weil DU, Tarok, uns angegriffen hast. DU mit deiner ewigen Anfeindung und dem Schmerz, den du an uns auslässt. Ich war genauso entsetzt über die tragischen Tode wie du, aber DU musstest ja einen gleich eine ganze Horde der deinigen auf die Besatzung von Karulzar losschicken.“

Kreatok schien wieder etwas sicherer zu werden, aber er wirkte definitiv nicht gesünder, als er ausspuckte: „Ich lag falsch. Es war nicht Nehals Schuld. Es war aber auch nicht meine. Es war deine Schuld! Was nach diesem tragischen Unfall geschehen ist... all diese Tode hast DU auf dem Gewissen!“

„DU HAST IN DEN FLAMMEN GESUNGEN UND GETANZT! IN DEINEN KÜHNSTEN TRÄUMEN SAH DAS NICHT NACH EINEM UNFALL AUS!.“

Kreatok hielt inne. Er setzte zu einer Antwort an: „Das... das war nicht... der Schild, er trägt Dunkelheit in sich... sobald ich ihn einmal abgelegt hatte, da...“

Der Meisterschmied blickte den Schild an, den er immer noch festhielt, und hielt ihn etwas zur Seite, als würde er ihn jeden Augenblick fallen lassen.

Und Gouri sah, wie Tarok seinen Hals aufrichtete und tiefer ein- und ausatmete. Sie erkannte den Plan des Drachen. Sobald Kreatok den Dunkelschild loslassen würde, würde Tarok dessen Leben wie eine Kerze auspusten. Und dann wären die Seiten im Krieg wieder im Ungleichgewicht. Die Schildzwerge würden ausradiert werden. Prip, welcher gerade verzweifelt versuchte, eine am Boden liegende Kameradin vom Schlachtfeld wegzuzerren, würde sterben, wenn nicht durch das Feuer, dann durch die Trolle, welche nach jedem Scharmützel im Grauen Gebirge die Verletzten zusammensammeln kamen.

Das durfte sie nicht zulassen.

Doch was konnte sie schon tun?

Und in diesem Augenblick schien die Zeit sich zu verlangsamen, und Gouri erkannte etwas Glitzerndes auf Taroks Rücken. Sie schaute genauer hin. Zwischen zwei Stacheln, wie sie zu Dutzenden aus seiner Wirbelsäule herausstachen, sass sicher eingeklemmt ein bläulich glänzender Schild.

Belenor hatte vorhin erwähnt, dass der Sternenschild an die Drachen gefallen war. Konnte es sein, dass Tarok ihn mit sich gebracht hatte, als letzter Schutz gegen das Feuer des Dunkelschilds?

Wenn der Dunkelschild die Verkörperung des Finsteren und Verdorbenen war, so war der Sternenschild die Verkörperung des Guten und Reinen. Wenn nur jemand an den Sternenschild gelangen könnte, so könnte man damit das Schicksal vielleicht wieder ins Lot wenden. Doch dafür müsste jemand zu Tarok, auf dessen Rücken klettern, und den Sternenschild stibitzen, ehe die riesige Echse etwas merkte. Das war eine Tat, die an Unmöglichkeit angrenzte. Und dazu kam noch, dass niemand in der Nähe war, dem Gouri diesen Plan überhaupt mitteilen könnte. Sie selbst war wahrlich nicht dafür geeignet, so unerfahren und schwach, wie sie war. Gouri war kurz vor dem Verzweifeln – da tauchte plötzlich vor ihrem inneren Auge das Gesicht eines freundlich glucksenden Trolls auf.\bigskip



Der Troll erwachte aus seinem steinernen Schlaf und gähnte ausgiebige. Langsam rappelte er sich auf und räkelte sich im Licht der Morgensonne.

Er hatte keinen Namen.

Die niederen Trolle benötigten keinen Namen.

Er hatte nur seinen Clan, und die Bedürfnisse seines Clans kamen vor seinen eigenen.

Seltsam, dass er etwas abseits des Clan-Lagers eingenickt war. Sein Magen grummelte, aber das war es nicht, was ihn geweckt hatte. Als er auf seine Hand blickte, lagen darin die Überreste eines verkohlten Schildes.

Ach ja, genau! Der Troll hatte letzte Nacht zwei wirklich nette Gesellen getroffen, die seine Brücke überquert hatten, und diese Gesellen hatten ihm die verkohlten Überreste eines mythischen Schilds geschenkt. Der Troll grinste selig. Wahre Freunde waren das gewesen. Nicht weggerannt waren sie, wie es sonst alle taten, sondern auf ihn zugegangen und mit ihm gespielt hatten sie. Er gluckste schon allein von der Erinnerung an das Treffen. Leider waren sie davongerannt, ehe es Zeit fürs Essen geworden war – sie hatten so lecker ausgesehen! Und leider hatten ihre mythischen Schildüberreste nichts Tolles drauf gehabt.

Aber das machte nichts. Schliesslich besass sein Clan noch immer den braunen Schild mit dem Händeschütteln drauf. Und das war ein wahrlich toller Schild, das hatte der Troll schon von weitem gespürt.

Da tauchte vor dem inneren Auge des Trolls das Gesicht des einen dieser Gesellen auf, ein grünes Zwerglein mit einer Kleidung aus Moos war das gewesen. Und diese Gestalt schwebte nun in seiner Vorstellung rum und grinste ihn an. Das war ja ungewöhnlich. Was noch ungewöhnlicher war: Die Gestalt gestikulierte in eine gewisse Richtung – weg von der Brücke, und hin zu seiner Clanhöhle.
  
Fröhlich folgte der Troll dem grinsenden Gesicht der grünlichen Gestalt in die Höhle seines Clans. Der Wächter-Troll guckte ihn seltsam an, aber unserem Troll war das egal. Es fühlte sich gut an, diesem Gesicht zu folgen. Der Clan-Häuptling guckte noch dümmlicher drein, als der Troll einfach so in seinen Schlafraum stapfte und zum dort an der Wand hängenden mythischen Schild trat, aber auch das war unserem Troll egal. Unser Troll fühlte sich sicher und geleitet. Das Lächeln der grünen Person vor seinem inneren Auge bestätigte ihn in seinem Tun. Ein letztes Mal blickte er ihr in die Augen, diesmal fragend. Das Gesicht nickte, und so griff der Troll nach dem Schild und hob ihn in die Höhe und –

– da fühlte der Troll seine Knie weich werden und nachgeben. Was für finstere Hexenkunst steckte in diesem Schild? Der Troll wurde von einem Augenblick zum nächsten immer schwächer und schwächer! Schon mochte er den verfluchten Schild nicht mehr halten, und der Schild kullerte davon und rollte seinem ihn immer noch mit grossen Augen anstarrenden Clan-Häuptling vor die Füsse.

Unser Troll wollte nach dem Schild greifen, doch seine Arme waren zu schwer, viel zu schwer, und plumpsten nutzlos auf den Höhlenboden. Dann tat es ihnen sein restlicher Köper nach.

Da lag der Troll nun zusammengesunken am Boden. Selbst das Atmen kam ihm wie eine riesige Anstrengung vor. Was war nur los hier? Das immer noch freundlich lächelnde Gesicht der grünlichen Gestalt verblasste.\bigskip



Gouri spürte die Veränderung kommen, ehe sie eintraf. Wie die Spannung, die vor einem Gewitter in der Luft liegt, war ihr Körper aufgeladen und bereitete sich auf die einströmende Kraft vor. Breitbeinig stand Gouri oben auf der Hügelkuppe und starrte in die Ferne, das Gesicht des freundlichen Trolls immer noch im Kopf behaltend. Da spannten sich plötzlich ihre Muskeln unverhältnismässig an, und ein Blitz aus reiner Energie zuckte durch ihren Körper, drang in jede Faser ihres Wesens und erfüllte sie von Kopf bis Fuss. Gouri schrie schmerzerfüllt auf.

Dann war es vorbei.

Aber alles hatte sich verändert.

Gouri fühlte sich so leicht, so leicht, wie sie sich noch nie in ihrem Leben gefühlt hatte, leicht wie eine Feder, und doch so standhaft wie das Graue Gebirge selbst. Ihr Atem ging rasch, und sie fühlte, als könnte sie Bäume ausreissen. Konnte sie wahrscheinlich tatsächlich.

Ein Hoch darauf, dass sie sich an den Aufenthaltsort des Bruderschilds erinnert hatte. Ein Hoch auf den netten Troll, der sie in einem so positiven Licht sah, dass er ihr seine Stärke geliehen hatte. In gewissem Sinne auch ein Hoch auf Kreatok und Nehal, die diesen so mächtigen Bruderschild geschaffen hatten. Gouri jubelte. Sie war bärenstark, nein, sogar trollstark! Jetzt galt es, den Sternenschild von Taroks Rücken zu erobern.

Gouri duckte sich, holte Anlauf, und tat einen SPRUNG, welcher sie meterweit in die Luft katapultierte. Windstösse trafen sie im Gesicht, und sie ruderte ungeschickt mit den Armen, ehe die grasbewachsene Erde unter ihr allzu schnell wieder auf sie zukam und sie unschön zu Boden krachte. Gouri kullerte den Hügel hinunter, wobei sie sich noch mehr Schürfwunden an den Armen und Beinen zuzog.

Ächzend richtete Gouri sich wieder auf und strich sich Matsch aus dem Gesicht.

Ein grossartiger Start. Aber positiv bleiben: Sie war immer noch stark wie ein Troll, einige Schürfwunden würden ihr da nichts ausmachen.

Gouri versuchte es erneut, diesmal etwas vorsichtiger. Selbst das Laufen war schwer, da ihr Körper jede Bewegung viel schneller vollzog, als sie es erwartete. Aber nach einigen Schritten hatte Gouri den Dreh soweit raus, damit sie ein etwas schnelleres Tempo einlegen konnte. Und so begann Gouri zu rennen.

Gute Güte, und wie sie begann, zu rennen! Die Umgebung um sie verzerrte sich zu grünen und braunen Farbschlieren, als Gouri einen weiteren Zahn zulegte. Schemen von verletzten und einander helfenden Zwergen huschten links und rechts an ihr vorbei, einige drehten sich erstaunt nach ihr um. Sie mochten viel Schauriges erlebt haben in den letzten Tagen, doch eine Agren mit der Stärke eines Trolls und der daraus resultierenden Geschwindigkeit sah man auch nicht alle Tage.

Gouri wagte einen kleinen Hopser, und schaffte es diesmal, nach der Landung auf den Beinen zu bleiben. Sie verlor kurzzeitig den Rhythmus, wurde langsamer, fand die Kontrolle wieder und beschleunigte auf ein noch höheres Tempo als vorher, stetig auf den riesigen schwarzen Drachen zu.

Kreatok und Tarok warfen sich weiterhin Anschuldigungen an den Kopf, doch Gouri kümmerte sich nicht mehr gross um den Inhalt. Kreatok hatte den Dunkelschild immer noch nachdenklich in der Hand und Tarok war immer noch komplett auf ihn fokussiert. Soweit, so gut.

Gouri wagte einen noch grösseren Satz, dann noch einen, und diesmal erreichte sie wieder eine Höhe von mehreren Metern. Das machte beinahe Spass!

So näherte sich Gouri hüpfend, immer höher hüpfend, dem riesigen Drachen. Sie sprang in hohem Bogen über den einem anderen Schildzwerg helfenden Prip, welcher sich aufrichtete und sie mit grossen Augen anstarrte. Von weit hinten hörte Gouri ihn etwas rufen, doch der Wind verzerrte es, sodass sie es nicht hören konnte. Kurz überlegte Gouri, stehenzubleiben, und gelang ins Stolpern. Dann rief sie sich ins Gedächtnis: Der Sternenschild. Sie musste zum Sternenschild, solange Kreatok den Dunkelschild noch nicht losgelassen hatte.

Nun war Gouri kaum mehr eine Drachenlänge von dem gigantischen Tarok und dem noch viel gigantischeren schwarz-silbernen Feuersturm Kreatoks entfernt. Die Hitze schlug ihr entgegen und brannte in ihren Schürfwunden. Ihre Schuhe hatte sie längst abgelaufen, und ihre Füsse würden die Tortur auch nicht mehr lange mitmachen. Aber weit war es nicht mehr, nur noch zwei, drei Schritte –

– da riss Kreatok mit einem inbrünstigen Schrei den Dunkelschild in die Höhe und der schwarze Flammensturm setzte sich in Bewegung, umströmte ihn und raste auf Tarok zu, welcher laut aufbrüllte und seinerseits einen feuerroten Flammenstrahl in Richtung des Schildzwergs sandte – und funkelten da etwa Sterne in Taroks Strahl?

Die Feuerwände krachten zusammen in einem Schauspiel, welches zugleich unglaublich furchterregend und wunderschön majestätisch aussah. Silberne Flammen trafen auf orange, schwarze auf rote, und die Feuer vermischten sich in Wirbeln und Strudeln, verglühten die letzten Reste des Grases zwischen dem Schildzwerg und dem Feuerdrachen und tauchten die Umgebung in ein gespenstisches rotes Licht. Silbern glitzernder Rauch stieg in die Höhe und verknüpfte sich mit den schneeweissen Wolken, über denen immer noch Drachen flogen und Lava auf die Drei Wasser niederregnen liessen.

Gouri war zu spät, um auszuweichen. Sie flog mitten im Sprung auf die vielfarbige Feuerwand zu, während die Hitze sich auf ein unerträgliches Mass steigerte.

Sollte sie noch versuchen, ihre Geschwindigkeit zu bremsen? Falls sie das überlebte, könnte Gouri sich zurückziehen und später einen erneuten Versuch machen, den Sternenschild auf Taroks Rücken zu erringen. Doch würde das Schicksal ihr eine zweite Chance geben?

Sie durfte den Troll nicht vergessen, dessen Stärke sie im Moment führte. Dieser Troll lag gerade hilflos zusammengesunken auf einem Höhlenboden nicht weit von hier entfernt, und brauchte dringend seine Kraft zurück. Sobald Gouri ihm diese zurückgegeben hatte, würde sie sie wohl nicht wiedererhalten können. Es hiess jetzt oder nie.

Gouri dachte an Prip und seine Zwerge, die sich selbstlos auf das Schlachtfeld gewagt hatten, um ihre Mitzwerge zu unterstützen. Ihre Bemühungen wären alle für nichts, wenn Tarok und Kreatok jetzt einen Feuersturm über das Schlachtfeld entfachen würden. Prip würde sterben. Belenor mit seinem prächtigen Schnurrbart würde sterben. Gouri selbst würde sterben, wenn sie nicht schleunigst Leine zog.

Und so, mitten im freien Fall, traf Gouri ihre Entscheidung.

Sie prallte auf dem Boden auf und sprang erneut hoch, viel zu schnell, viel zu weit, viel zu nahe an den tobenden Kampf der Feuer. Ungehindert schoss Gouris Körper direkt auf das Flammenmeer zu. Sie schloss die Augen und bereitete sich auf ihr Ende vor.

Ein kühler Luftstoss traf Gouri im Gesicht und fegte die Hitze davon. Ungläubig öffnete sie die Augen, gerade rechtzeitig für den nächsten Sprung. Tarok sass nicht mehr auf der Erde, sondern hatte seine riesigen Flügel aufgespannt und schwang sie würdevoll, sich langsam in die Lüfte erhebend.

Das schwarze Feuer konnte ihm gefährlich werden. Der Feuerdrache wollte sich in Sicherheit begeben. Und durch das Bewegen seiner Flügel sorgte er für frische Luftströme, die Gouri ganz nebenbei die rettende Kühle verschafften. Sie dankte der Mutter des Steins, und legte die letzten paar Meter in einem mächtigen Sprung zurück.

Im Nachhinein hätte sie besser zielen können. Aber Gouri war sowohl geistig als auch körperlich an ihrem Ende, und hatte sich nur auf den Sternenschild fokussiert, darauf, dass der Schild gleich aus ihrer Reichweite verschwinden würde und dass sie ihn vorher erreichen musste. So hatte sie zum letzten Sprung angesetzt und sich weit in die Lüfte erhoben, an Taroks Beinen, Brust und selbst über seine Kopfhöhe hinweg, direkt auf seinen Rücken zu, wo zwischen zwei spitzen Stacheln weiterhin der Sternenschild eingeklemmt war. Aus der Nähe glitzerte er noch viel schöner.

Gouri lächelte.

Dann prallte sie auf den Rücken des Tarok und wurde gleich wieder meterhoch in die Luft geschleudert, knallte an einen seiner mächtigen Flügel, der sich gerade am Heben war, und rutschte daran herunter, bis sie am unteren Flügelansatz sass, beinahe nicht mehr bei Bewusstsein.

Gouri, die Agren, war auf Tarok gelandet.

Hatte er sie bemerkt? Ein rascher Blick verriet ihr, dass Tarok immer noch viel zu sehr damit beschäftigt war, einen Feuerstrahl nach dem anderen auf Kreatok abzusenden, als dass er sich darum kümmern konnte, was für mickrige Wesen da auf seinem Rücken landeten.

Gouri hatte Mühe, die Situation zu begreifen. Sie lag auf dem Rücken von Tarok. Sie lag auf dem Rücken eines vom abhebenden Drachen. Man könnte argumentieren, dass sie einen Drachen ritt – sie, eine einfache Agren aus dem Grauen Gebirge!

Durch den Rauch, welcher vom Drachenfeuer und vom Feuer des Dunkelschilds hoch in die Lüfte stieb, konnte Gouri den Boden nur noch schwer erkennen, aber es war klar, dass er sich von ihr entfernte. Mit mächtigen Schwüngen seiner breiten Flügel schraubte Tarok sich weiter in den Himmel hervor, weiter weg vom riesigen silbernen Feuer, welches Kreatok beschworen hatte. Gouris Höhenangst schlug rasant Alarm, also wandte sie sich vom Blick nach unten ab und drehte sich mühsam auf den Rücken.

Der Blick nach oben löste in ihr fast noch mehr Furcht aus. Über ihr kreisten mindestens fünf Drachen, welche aufgehört hatten, Lava auf die Drei Wasser regnen zu lassen, und stattdessen Taroks Kampf mit Kreatok aufmerksam beobachteten. Konnten die Drachen sie vielleicht erkennen?

In diesem Moment lief Gouri ein weiterer Schauer über den Rücken, und sie sackte etwas zusammen. Ihr Muskeln entspannten sich kurz und verspannten sich gleich wieder, als eine Welle reinen Schmerzes ihren Körper ergriff.

Gouri brauchte einige Augenblicke, um zu erkennen, dass es sich dabei nicht um einen telepathischen Angriff eines Drachen handelte. Nein, es war die Kraft, die sie vom Troll erhalten hatte. Die Stärke des Trolls verliess ihren Körper und ihre eigene Verletzlichkeit trat wieder zum Vorschein. Gouri schaute auf ihre Hände, welche stark zitterten und ihr nicht mehr zu gehorchen schienen. Ihr Blickfeld verengte sich und die Welt wurde dunkel. Blind warf Gouri sich nach vorne, auf den mächtigen Sternenschild zu, welcher nur wenige Schritte von ihr entfernt zwischen den Rückenstacheln des Tarok steckte.

Sie verlor das Bewusstsein nicht. Sobald ihr Körper den Schild berührte, ertönte ein leises metallisches Summen, und der Schleier vor ihren Augen lichtete sich. Die Schmerzen liessen nach. Der Sturm von Taroks Flügeln wurde leiser und leiser. Und Gouri fühlte nicht mehr Taroks erhitze Schuppen unter ihren Händen, sondern festen, kalten Stein.

Verwirrt blickte Gouri sich um. Sie erblickte dunkle Äste und Blätter auf allen Seiten, weiter vorne sogar einen riesigen Stamm. Offenbar befand sie sich unter einem steinernen Mammutbaum, höher als alle Burgen, die sie in ihrem Leben je gesehen hatte, so gross, dass seine Blätter den Himmel verdeckten. Schimmernde Edelsteine überzogen den steinernen Baum, und gleissendes bläuliches Licht strahlten diese aus. Dieses Licht umgab Gouri, und dieses Licht spiegelte sich im glänzenden Sternenschild, den sie immer noch in der Hand hielt.

Tatsächlich: Sie besass den Sternenschild! Eines der mächtigsten Artefakte, von denen sie aus zahlreichen Legenden und Sagen gehört hatte. Sagen und Legenden, denen sie als Kind begeistert gelauscht und die sie in sich aufgesogen hatte. Und nun war sie selbst mitten in einer Sage.

Weiter vorne, nahe beim Stamm, rührte sich etwas. Eine seltsame Kreatur mit Hornklauen anstelle von Händen löste sich aus einem Loch im Stamm des riesigen Steinbaums und krabbelte geschickt über den kalten Stein auf Gouri zu. Die Kreatur erreichte sie jedoch nie, denn in jenem Augenblick schoss eine gewaltige Gestalt an ihr vorbei und landete zwischen der Kreatur und Gouri.

Es war ein Drache, wie konnte es auch anders sein.

Gouri erkannte erst jetzt die übrigen fliegenden Drachen, wie sie weit über ihr den Steinbaum umkreisten. Viele von ihnen schenkten ihr keine Beachtung, aber manche verdrehten die langen Schlangenhälse, um sie zu erkennen. Einer der Drachen, welcher besonders schnell den Baum umkreiste, erinnerte Gouri ungut an Tarok – und wenn sie plötzlich an diesem Ort sein konnte, so war dieser Drache vielleicht auch Tarok selbst.

Als es vor ihr schnaubte, wandte Gouri sich rasch wieder dem Drachen zu, welcher sich auf dem Boden niedergelassen hatte und sie aufmerksam studierte:

„Der Sternenschild?“, fragte eine tiefe, aber nicht unfreundliche Stimme in ihrem Kopf, „Ich hätte nicht gedacht, dass ich ihn so bald schon wiedersehen würde. Du brauchst dich nicht zu fürchten, kleine Agren. Solange dein Geist hier in Krahal weilt, können wir dir nichts tun. Also setzen wir uns hin und reden über diese... Angelegenheit.“

Der Drache strahlte eine gewisse Ruhe und Gelassenheit aus, und Gouri spürte, dass er die Wahrheit sagte.

Tarok (falls es denn Tarok war) hatte sich Fleckchen Stein, auf dem Gouri und der fremde Drache sassen, nun ebenfalls angenähert, und schwebte elegant in der Luft, jedoch ohne sich niederzulassen. Er überragte den ruhigen Drachen um mindestens das dreifache.

„DU!“, zischte Tarok zu Gouri, „Du hast meinen Sohn Sagrak auf dem Gewissen!“ Doch zu einem physischen Angriff liess Tarok sich nicht hinreissen.

Der Sternenschild summte plötzlich auf, und Gouri spürte eine Welle des Schmerzes von Tarok ausgehen. Erlaubte der Schild ihr, die Emotionen der Drachen in Krahal zu teilen? Sie spürte einen vertrauten Geist aus der Höhe, und als sie ihren Hals verdrehte, erkannte sie Sagrak, wie er weit über ihr auf einem steinernen Ast sass und sie aufmerksam bespitzelte. Waren die meisten Drachen, die sie hier sah, nur Abbilder bereits gestorbener Drachen? Waren es ihre Geister? Gouri hatte Geschichten von Krahal gehört, doch jetzt, wo sie den Ort mit eigenen Augen sah, war es einfach absolut überwältigend. Was tat ihr Körper momentan in der realen Welt? Und was Tarok?

Als hätte er ihre Gedanken erraten, antwortete der ruhige Drache nun:

„Keine Sorge, du stürzt nicht gerade vom Himmel. Tarok vermag es, seine Flügel zu schwingen, auch wenn sein Geist gerade in Krahal ruht.“

Mit einer gekrümmten Kralle deutete der ruhige Drache in die Höhe. Was Gouri zunächst für den vom riesigen Steinbaum verdeckten Himmel gehalten hatte, wirkte nun wie eine Höhlendecke. Doch das konnte natürlich nicht sein, schliesslich hatte eine so grosse Höhle unmöglich unter der Erde Platz.

Wichtiger war jedoch, was sich an dieser möglicherweise-Höhlendecke abspielte: Darauf erkannte Gouri eine sich bewegende Wandmalerei, ein Abbild ihrer selbst, wie sie auf dem Rücken Taroks lag, den Sternenschild fest umklammert. Dieses Abbild von Tarok bewegte sich nicht gross, aber seine Flügel schlugen ununterbrochen. Durch die Decke konnte Gouri erblicken, was sich soeben in der realen Welt abspielte.

Der ruhige Drache seufzte, und sprach: „Solange du den Sternenschild trägst, gebietest du darüber, wer Krahal sonst noch zu betreten vermag. Ich bitte dich darum inbrünstig: Wärst du so lieb, Kreatok hierhin zu bringen?“

Gouri atmete tief ein und aus und streckte ihre Hand schützend aus: „Ich... ich brauche kurz einen Moment. Das... das ist alles etwas zu viel.“

Der ruhige Drache seufzte wieder und liess sich zu Boden sinken, die Vorderpfoten elegant gekreuzt: „Du hast alle Zeit der Welt. Naja, ‚alle Zeit‘ ist vielleicht etwas übertrieben. Aber jedenfalls genügend Zeit.“

Tarok, welcher vorhin noch unruhig über ihnen beiden hin- und hergetigert war, atmete tief ein und aus und liess sich dann auch zu Boden fallen. Die Erde zitterte unter Gouris Füssen und sie fühlte sich plötzlich ganz klein, direkt vor der Schnauzen zweier Drachen, die sie anstarrten.

Der ruhige Drache wandte sich Tarok zu und fragte vorsichtig: „Bist du sicher, dass hilfreich ist, wenn du dich so nahe befindest...“

Tarok drehte seinen Kopf dem ruhigen Drachen entgegen und öffnete seinen Mund leicht, ein furchterregendes Knurren ausstossend, woraufhin der viel kleinere ruhige Drache unterwürfig den Kopf senkte und Tarok gewähren liess. Hitzewellen entwichen Taroks Maul und trafen Gouri im Gesicht, die sich abzuwenden hatte. Erneut fragte sich, wie sicher sie hier in Krahal wirklich war. Was würde geschehen, wenn dieser Tarok – sein Geist – sich entscheiden würde, sie zu beissen? Würde sie in der echten Welt aufwachen oder konnte sie hier drin sterben?

Plötzlich glomm der Sternenschild in einem beruhigenden blauen Licht auf, und Gouri überkam die Gewissheit, dass sie im Moment nicht gefährdet war. Sie war sich absolut sicher, dass Tarok sie nicht verletzen konnte. Der Schild würde sie beschützen. Der Schild hatte Hoffnung in ihr geweckt.

Der ruhige Drache blickte vom wütenden Tarok auf und fast so etwas wie Stolz überkam ihn, als er den leuchtenden Schild ansah und sprach: „Ach... ist er nicht wunderschön? Von den vieren war er mir stets der liebste.“

 Gouri sah ihn überrascht an. Der ruhige Drache sprach weiter: „Verzeih mir, in diesem ganzen Klamauk hatte ich mich noch gar nicht vorgestellt: Ich bin Nehal. Ich war derjenige, der mit Kreatoks Hilfe die vier mächtigen Schilde hergestellt hatte. Und auch wenn Kreatok jeweils den Silberschild benutzt hatte, um nach Krahal zu gelangen, denke ich, dass mit geschickter Führung auch eine Trägerin des Sternenschilds seinen Geist hierhin bringen könnte. Kannst du das für uns tun, edle Dame?“

Gouri starrte Nehal an. Der Erschaffer der mächtigen Schilde! Kurz hatte sie das Verlangen, sich zu verbeugen – sie hatte bei Prip abgeschaut, wie das ging – doch dann liess sie es sein. Sie hatte einhundert Fragen und nochmals einhundert mehr, doch das war nicht der Zeitpunkt, erst recht nicht vor dem grimmig dreinschauenden Tarok. Somit begnügte sich Gouri damit, ein leises „Ich bin Gouri“ zu stammeln und dann zu fragen: „Wie... wie kommt es, dass du... dass Sie... hier...“

„Ah, wir Drachen sind Seelen der Erde und des Feuers, wir sind nicht einfach so auszulöschen, indem man unsere Körper vernichtet“, lächelte Nehal, „Solange noch einer einziger von uns übrig ist, leben wir in dessen Kopf weiter. Krahal besteht weiter, solange noch einer von uns bei klarem Verstand ist“ – bei diesen Worten schaute Nehal warnend auf Tarok – „und das wird wohl noch jahrtausendelang der Fall sein.“

„KOMMT KREATOK JETZT ODER NICHT?“, warf Tarok ungeduldig ein.

Gouri und Nehal verstummten umgehend.

„Wenn ich Kreatok hierhin bringe, könnt ihr ihm nichts tun?“, fragte Gouri nochmals nach.

„Ja, Träger der mächtigen Schilde sind hier sicher. Wir wollen nur reden.“, antwortete Nehal, und Gouri spürte abermals durch den Sternenschild in ihrer Hand, dass Nehal dies als die reine Wahrheit ansah.

So weit, so gut. Nun musste Gouri nur noch Kreatok nach Krahal holen. Sie hatte zwar keine Ahnung, wie sie das anstellen sollte, aber sie hatte Vertrauen in die mächtigen Schilde – bislang hatte sie ja auch nicht mehr Ahnung gehabt, wie die Schilde genau einzusetzen waren.

Und prompt summte der Sternenschild metallisch, glomm leuchtend auf und in Gouris Nähe verdichteten sich die Nebelschwaden. Heraus trat Kreatok, der Meisterschmied, den Dunkelschild immer noch in der Hand. Das lief ja wie am Schnürchen.

Im Gegensatz zu Gouri wirkte Kreatok nicht überrascht, in Krahal zu stehen. Sein gefasster, fast gelangweilter Gesichtsausdruck deutete darauf hin, dass er bereits einmal hier gewesen war. Das machte auch Sinn, schliesslich hatten er und Nehal die mächtigen Schilde geschaffen und wussten ziemlich sicher über alle Tricks Bescheid, die diese auf Lager hatten.

Da er nicht über den riesigen Baum inmitten Krahals und die Unmenge an herumfliegenden Drachen staunen musste, konnte Kreatok seine Aufmerksamkeit direkt auf Tarok richten. Tarok rührte sich nicht vom Fleck, aber sein Schwanz zuckte immer wieder und seine Augen verengten sich – wenn Gouri die Körpersprache des riesigen Reptils richtig deutete, stellte er sich gerade vor, wie nett es wäre, den kleinen Zwerg zu rösten.

Kreatoks Mine verfinsterte sich. Er nickte Gouri mit einem kalten Gesichtsausdruck kurz zu und setze zu einem Spruch an – wahrscheinlich wollte er schnippisch fragen, warum sie ihn nach Krahal gebracht hatte.

Dann fiel Kreatoks Blick auf Nehal, und er blieb mit halb geöffnetem Mund stehen. Beinahe sofort begannen seine Augen zu glänzen. Tiefe Furchen erschienen in seiner gerunzelten Stirn und sein Kinn zitterte, sodass sein ganzer Bart zu wackeln begann. Ja, Gouri fiel auf, dass sein Bart, welcher in der Welt da draussen schon längst versengt worden war, hier in Krahal wieder in voller Pracht vorhanden war. Was für ein mysteriöser Ort.

Nehal erwiderte Kreatoks Blick aus rot glühenden Augen, aber im Gegensatz zu Taroks leuchteten sie nicht bedrohlich, sondern warm und einladend, wie der Schein eines warmen Lagerfeuers. Nehal rief seinen Mörder beim Namen, in einer Tonlage, in welcher man sonst einen alten Freund begrüsste: „Kreatok...“

Eine einzelne Träne löste sich aus Kreatoks Augenwinkel und rollte seine Wange entlang. Sein Gesichtsausdruck zeigte nicht mehr den verwirrten entflohenen Gefangenen, als den Gouri ihn kennengelernt hatte. Und er war auch nicht mehr der kühle, spöttische Zwerg, der sich Tarok in den Weg gestellt hatte. Kreatok wirkte jetzt wie jemand, der jahrelang in der Dunkelheit gesessen hatte und nun einen Strahl Sonnenlichts erblickte. Wie jemand, dessen Kind an einer schweren Krankheit gelitten hatte und der nun von dessen Heilung erfuhr. Wie jemand, der seinen besten Freund im Affekt umgebracht hatte und ihn nun wieder erblickte, nicht wütend oder enttäuscht, sondern einfach... war das Vergebung, das Nehals Gesicht ihm versprach? Konnte ihm überhaupt vergeben werden?

Gouri spürte abermals durch den Sternenschild, dass Kreatok der Lage noch nicht traute. Er selbst hatte sich nie verziehen und hatte in den Jahren der Einsamkeit in den Kerkern der Drei Wasser gelitten, als er sich immer und immer wieder vorgestellt hatte, was Nehal nur von ihm denken musste. Kreatok hatte beim Schmieden des Dunkelschilds, von Nehal unbemerkt, machthungrig eine Dunkle Quelle angezapft, und diese hatte ihn erfüllt, begleitet, geleitet, bis er nicht mehr sie kontrollierte, sondern sie ihn. Erst als er Nehals Schreie gehört hatte, war Kreatok wieder aufgewacht, und hatte den Dunkelschild von sich gestossen. Doch zu diesem Zeitpunkt war es bereits zu spät gewesen. Nehals Körper war vernichtet und Kreatok hatte den Krieg zwischen Zwergen und Drachen ausgelöst.

Ein Bild blitzte vor Gouris innerem Auge auf: Die rauchenden Überreste Nehals, wie sie umgeben von silbernem Feuer in einer Höhle tief unter Cavern brannten. Sie wusste, dass dieses Bild in Kreatoks Erinnerung festgebrannt war, und wie schon seit vielen Jahren kämpften in seinem Innern zwei Seiten.

Die eine Seite Kreatoks wollte den Dunkelschild wieder von sich werfen, so weit weg, wie nur irgend möglich, und zu Nehal rennen, ihn umarmen, ihm sagen, dass es ihm so unglaublich leid tat. Doch Worte konnten dem, was er getan hatte, nicht gerecht werden.

Die andere Seite Kreatoks konnte den Schild nicht loslassen. Kreatok war nach so langer Zeit zum ersten Mal wieder mit einer seiner Kreationen vereint, und sie war seine Lebensversicherung. Ohne Schild würden die Schildzwerge ihn wieder gefangen nehmen, oder schlimmer, Tarok ihn direkt vernichten. Nein, schon die ganze Angelegenheit hier musste ein Trick sein. Nehal konnte in ihm nur den Mörder sehen, den er selbst in sich sah. Sobald er den Schild losliess, würden sich die beiden Echsen auf ihn stürzen und seinen Geist zerreissen. Er brauchte den Dunkelschild.

Und so blieb Kreatok stehen, den Dunkelschild immer noch fest gepackt, und starrte Nehal entgegen, immer mehr Tränen in den Augen.

Auch von Nehal gingen jetzt Wellen der Trauer aus, als er seinen alten Freund vor sich anblickte, wie er zitterte und den Dunkelschild umklammerte. Erneut sprach er seinen Namen, diesmal mit einem leichten Anschwung von Tadel in der Stimme: „Kreatok... du brauchst diesen Schild doch nicht.“

Immer noch strahlte Nehal Vergebung und Freude aus, als er hinzufügte: „Komm, alter Freund. Du musst nicht mehr leiden.“

Nehal hatte seine Emotionen nicht mehr vollständig unter Kontrolle und für einen kurzen Augenblick fühlte Gouri, wie viele Dinge Nehal noch sagen wollte und nicht auszudrücken vermochte, weil Worte nicht genug waren. Es muss dir nicht leidtun, Kreatok. Du bist nicht das Böse hier. Es wird alles in Ordnung sein.

Kreatok spürte diese Dinge auch, und er war so nahe davor, den Dunkelschild loszulassen. Der Sternenschild in Gouris Händen summte ein letztes Mal auf und strahlte heller denn je zuvor, als die Agren alle ihre Zuversicht auf den Schildzwerg lenkte, wie er da vor ihr, Nehal und Tarok stand und mit sich selbst haderte.

Etwas in Kreatok brach und er murmelte: „Was soll’s.“

Dann öffnete er seinen Griff und der Dunkelschild entschlüpfte ihm, zu Nebel verpuffend, sodass kein Abbild des Schildes mehr in Krahal blieb. Gouri fiel auf, dass Kreatok auch hier in Krahal noch die Überreste goldener Ketten trug. Sie brauchte nicht nach oben in die echte Welt zu blicken, um zu wissen, dass der wahre Kreatok soeben den Dunkelschild losgelassen hatte.

Die Schlacht um die Drei Wasser war vorüber.

Nehal lachte fröhlich auf.

Kreatok schluchzte auf und rannte auf Nehal zu, die Arme zu einer Umarmung ausgestreckt.

Er sollte nie dort ankommen.

Tarok, der sich die letzten Minuten unnatürlich still verhalten hatte, bäumte sich brüllend auf und spie einen mächtigen Feuerstrahl auf den noch rennenden Kreatok, welcher innerhalb eines Augenblicks Krahal verliess und zu Nebel verpuffte. Erneut musste Gouri nicht nach oben in die reale Welt sehen, um zu wissen, dass Tarok soeben das Feuer auf Kreatok eröffnet hatte und Kreatoks Körper keine Chance gehabt hatte.

Gouri schrie auf: „NEIN!“

Nehal blickte wie versteinert auf Tarok. Hatte er gewusst, dass dies Taroks Plan gewesen war? Hatte er dies etwa unterstützt?!

Sie sollte die Antwort nie erfahren, denn in diesem Augenblick schrie Tarok: „DAS IST ES, WAS IHR VERRÄTER VERDIENT!“

Und der riesige Tarok stürzte auf Nehal zu, packte den kleineren Drachen im Sprung und hob schleuderte ihn in die Höhe. Nehal prallte überrumpelt und ungelenk auf den Boden, wo er mit einem verblüfften Gesichtsausdruck so rasch zu Nebel verpuffte, wie es Kreatoks Geist soeben getan hatte. Gouri war geschockt. Konnten die Überreste der Drachen in Krahal immer noch vernichtet werden?

Kreatok und Nehal. Die Erschaffer der vier mächtigen Schilde. Figuren aus Legenden, die Gouri als kleine Agren geliebt hatte. Und nun, so kurz nachdem sie erfahren hatte, dass die beiden noch existierten, hatte Tarok ihnen ein gemeinsames Ende bereitet.

Tarok zuckte erregt und wandte sich Gouri zu, Wut und Rache immer noch sein Denken beherrschend. Gouri packte den Sternenschild fester und erwiderte Taroks bösartigen Blick. Die Sache durfte nicht noch weiter eskalieren. Sie war bereit, für das Ende dieses Kriegs zu kämpfen.\bigskip



Prip schwitzte Blut und Wasser. Eben noch hatte er gesehen, wie Gouri mit überzwerglicher Geschwindigkeit an ihm vorbeigehüpft war, dem Flammeninferno vor ihm entgegen, und sich auf den Rücken des Tarok geschwungen hatte. Was bei allen sieben Feuern der Tiefminen war ihr bloss in den Sinn gekommen? Und woher hatte sie diese Stärke errungen?

Prip hatte sich wieder seinen verwundeten und gefallenen Kameraden zuwenden müssen, doch als er die nächste Schildzwergin aus ihrer Rüstung befreit hatte und in eine möglichst sichere Zone abseits der Festung zerrte, fiel ihm die Begegnung mit Troll und dem Bruderschild wieder ein. Der Troll musste tatsächlich seine riesige Stärke an die Agren übergeben haben. Prip wurde warm ums Herz... war das Stolz?

Gleich danach hüpfte ihm das Herz in die Hose, als ihm einfiel, dass Gouri, kräftig wie ein Troll oder auch nicht, soeben auf einem abhebenden Drachen gelandet war. Verschreckt schaute er sich nach dem Kampf der Giganten um.

Das war ja seltsam! Die Feuer, sowohl das silberne des Dunkelschilds als auch das rote des Feuerdrachen, waren erloschen. Kreatok stand immer noch vor den Wänden der Drei Wasser und schien sich nicht zu bewegen. Tarok hing in der Luft und schlug mit seinen riesigen Flügeln, dass die umliegenden Bäume schwankten, doch abgesehen davon bewegte sich der Drache auch nicht. Und Gouri konnte Prip gar nicht mehr erkennen.

Lieferten sich Kreatok und Tarok ein geistiges Gefecht?

Da! Tarok öffnete seinen Schlund und liess Feuer auf Kreatok herunterregnen! Der Körper des Zwergs wurde innert Sekunden verschlungen und zu Asche. Prip knirschte mit den Zähnen. Klar, Kreatok selbst würde niemand vermissen. Aber ohne ihn und den Dunkelschild wären die Drei Wasser endgültig verloren.

Doch Tarok schlug weiterhin bloss mit den Flügeln und blieb ansonsten regungslos. Lieferte er sich einen weiteren geistigen Kampf? Das konnte doch nur mit Gouri sein? Wo befand sich die Agren? Ging es ihr gut?

Prip trotte hinüber zu Belenor, welcher sich auf einem behelfsmässigen Lager von Kreatoks Giftpfeil erholte. Dort griff Prip nach dem Fernrohr, welches Belenor an seinem Gürtel trug. Auch eine Erfindung, welche Kreatok perfektioniert hatte. Prip versuchte nicht, daran zu denken, als er das Fernrohr auf Tarok richtete.

Da! Auf dem Rücken von Tarok, eingekeilt zwischen zwei spitzen Stacheln, erkannte er etwas Glitzerndes. War das etwa... der Sternenschild? Und da unter dem Schild, lag da nicht... Gouri! Die Agren hatte den hell leuchtenden Schild umklammert und lag so auf Taroks Rücken, sich ebenfalls nicht weiter bewegend. Sie trat tatsächlich im Geist gegen den Feuerdrachen an. Prip wischte sich den Staub aus dem Gesicht und verfluchte sich dafür, dass er nichts ausrichten konnte.

In diesem Moment erlosch das Glitzern des Schildes. Tarok richtete sich in der Luft auf und stiess ein grauenvolles Geheul aus, welches Prip durch Mark und Bein ging. Dann schwang er sich weiter in die Lüfte, immer noch kreischend, und drehte sich einmal um die eigene Achse. Prip saugte tief Luft ein, als er sah, wie Gouris winziger Körper den Rücken des Drachen verliess und dicht gefolgt vom Sternenschild in den Himmel geschleudert wurde.

Tarok, der furchteinflössende riesige Feuerdrache, schrie immer noch und wedelte mit seinen gewaltigen Tatzen, als wolle er einen unsichtbaren Feind abzuwehren. Das konnte nur etwas heissen: Der geistige Kampf war nicht zu seinen Gunsten ausgegangen! Die übrigen Drachen, die immer noch über der Festung kreisten, stiegen mit ein in Taroks Geheul und zogen mit ihm ab.

Jubelschreie von Seiten der Schildzwerge wurden laut und auch in Prips Herz machte sich Erleichterung breit, als er erkannte, dass die Dachenbrigade die Flucht ergriffen hatte – für den Moment jedenfalls. Nur der versteinerte Drache neben den Drei Wassern war noch übrig.

Doch Prip freute sich nicht so sehr wie die anderen. Seine Gedanken schweiften rasch zurück zu Gouri, und so rannte er los, der Stelle entgegen, wo ihr Körper gelandet haben könnte.

Bereits von weitem konnte er den Sternenschild glänzen sehen, trotz der Menge an Schlamm und Dreck, die ihn umgab. Und darunter lag eine Gestalt... Prip versuchte, seinen Schritt zu beschleunigen, doch die Verletzungen, die er sich im Kampf gegen Sagrak zugezogen hatte, zeigten sich umso deutlich. Mehr humpelnd als rennend legte er die letzten Meter zur Einschlagstelle des Sternenschilds zurück. Die ganze Zeit hatte er sich eingeredet, dass niemand einen Sturz von dieser Höhe überlebte, auch nicht jemand mit der Stärke eines Trolls – falls Gouri diese überhaupt noch besessen hatte. Und doch hatte er die Hoffnung noch nicht aufgegeben, als er den Sternenschild erreichte, sich neben ihm zu Boden fallen liess und den leuchtenden Schild ungelenk von Gouris Körper zerrte.

Da lag Gouri nun, in einem Krater vor der Festung, ihre eigentlich grüne Mooskleidung braun vom Dreck und rot vom Blut, das Gesicht nach unten, die verzottelten Haare grösstenteils versengt. Sie hatte Tarok irgendwie verschreckt und damit die Überreste der Drei Wasser gerettet. Das Zwergenvolk würde Zeit haben, sich zu erholen und sich zu wappnen für den nächsten Angriff. Sie besassen noch drei mächtige Schilde, darunter den Sternenschild!

Gouri hatte sich geopfert, um einem ihr fremden Volk Hoffnung zu schenken.

Was für eine heldenhafte Tat. Darüber würden eines Tages Geschichten geschrieben werden.

Prip kniete sich neben Gouri nieder. Er machte sich keine Illusionen. Sie atmete nicht mehr.

Prip schluckte tief, einen dicken Kloss im Hals. Prip bereite sich darauf vor, den Letzten Segen zu sprechen. Er hustete sich etwas Dreck aus der Lunge, und griff nach Gouris Schultern, um den Körper der Agren umzudrehen.

Prip konnte es kaum glauben: Gouris Gesicht war verkratzt und verschmutzt, aber wie durch ein Wunder grösstenteils unversehrt. Dasselbe... ja, dasselbe galt auch für ihren restlichen Körper. Keinerlei Knochenbrüche? Konnte das sein?

Ein metallisches Summen lenkte Prips Blick auf den Sternenschild, welcher immer noch neben ihm lag. Er leuchtete hell und im Schein dieses Lichtes konnte Prip erkennen, wie sich langsam jeder Schnitt, jede Prellung, jeder Schürfung auf Gouris Gesicht auflöste und Farbe darin zurückkehrte. Ein angenehmes Prickeln unter seiner Haut verriet Prip, dass dasselbe mit ihm geschah. Dann plötzlich brach das Summen des Schildes ab und ein stechender Schmerz zeigte Prip, dass er längst nicht vollständig geheilt worden war. Der Sternenschild hatte seine Kraft für heute wohl endgültig aufgebraucht.

Gouri sog tief Luft ein und verschluckte sich ganz grässlich. Eine Hustenattacke überkam sie und Prip vergass alle eigenen Wehwehchen.

Dann, endlich, schlug Gouri die Augen auf und lachte beinahe auf, als sie Prips besorgten Gesichtsausdruck sah.

In dieses glockenhelle Lachen stimmte Prip mit ein.\bigskip



Viel gab es zu tun in den nächsten Tagen. Die Drei Wasser waren beinahe komplett zerstört worden und viele gute Schildzwerge verletzt oder gar gestorben. Die Drachen konnten jeden Augenblick zurückkehren. Und so wurde die gesamte restliche Zwergenarmee unter dem Grauen Gebirge innert Kürze mobilisiert, um einen Rückzug der Besetzung der Drei Wasser unter die Erde, in die riesigen Gänge und Höhlen Caverns, zu vollbringen.

Dort tüftelte man bereits an trickreichen Fallen und Geschossen, welche auch von unter der Erde aus Drachen erlegen könnten. Denn der Unterirdische Krieg zwischen den Zwergen und Drachen war noch lange nicht zu Ende, er hatte vielmehr erst begonnen. Und sowohl die Trolle als auch die Krahder aus dem Süden würden etwas zu sagen haben darin.

Der Sternen- und Dunkelschild wurden zurück nach Cavern gebracht. Beide würden in den dunklen Tagen, die erst noch kommen sollten, wieder an die Drachen fallen. Einzig der Silberschild, in welchem kein lebender Zwerg oder Drache einen Zweck sah, verblieb während den gesamten restlichen Unruhezeiten in seiner Waffenkammer.

Prip und Gouri bekamen davon nur noch wenig mit. Sie halfen bei Bergungen und Aufräumaktionen mit, so gut sie konnten, und kamen sich in dieser Zeit näher. Am ersten Abend nach Kreatoks Tod versammelten sie sich neben den Massengräbern, um dem Meisterschmied und seinem Drachen zu gedenken. Sie waren wahrscheinlich die einzigen, die um die beiden trauerten.

Als der Abzug der Zwerge aus den Drei Wassern glücklich verlaufen war, verliess Prip das Zwergenvolk, um Gouri auf ihrem Heimweg zu begleiten. Der Rückweg verging ebenso wie die Aufräumarbeiten überraschend ereignislos. Keine Trolle warteten bei der Brücke und keine Drachen versuchten, die beiden Wanderer zu rösten.

Gouri war besonders nervös, denn sie hatte in all dem Chaos keinen Falken ausfindig machen können, um ihrer Familie zu berichten, dass bei ihr den Umständen entsprechend alles in Ordnung war. Die übrigen Agren mussten annehmen, dass Gouri längst bei einem Drachenangriff umgekommen war. Doch als sie dann endlich ihre Heimathöhle erreichte... nun, ich will nicht vorweggreifen. Lest einfach selbst:\bigskip



Alle waren versammelt in der Höhle: Gouris Eltern, ihre Geschwister, sogar der Alte Grone, der unlängst in den Seniorenkreis der Agren aufgenommen worden war.

Ihre Blicke landeten zunächst auf Gouri, als diese durch den Höhleneingang trat und sich schüchtern umsah. Ein überraschtes Raunen lief durch die Menge.

Dann wanderten die Blicke zu Prip, der sich in allerbester Heldenmanier (samt vergoldeter Armschiene und elegantem Gehstock) in eine Pose warf. Das in die Agrenhöhle scheinende Sonnenlicht liess seine Rüstung glorreich erglänzen. Das Murmeln der Agren nahm an Lautstärke zu.

Dann kehrten die Blicke der versammelten Agrengesellschaft zurück zu Gouri.

Ihr kleiner Bruder Krul war der erste, der sich rührte. Mit einem glücklichen Glucksen warf er die Moospuppe von sich, die er gehalten hatte, rannte auf Gouri zu und umarmte das Bein seiner grossen Schwester.

Damit war der Bann gebrochen. Ein Lächeln (vielerorts mit Freudentränen garniert) bildete sich auf den Gesichtern ihrer Eltern und Geschwister, und auch sie rannten nach vorne und umarmten ihr zurückgekehrtes Familienmitglied. Die restlichen Agren grinsten ebenfalls und kamen nach vorne. Einige musterten Prip noch argwöhnisch, doch die meisten kümmerten sich primär um Gouri und umkreisten sie, fragten sie, wo sie gewesen war, was sie mit ihrem Besitztum tun sollten, das sie schon auf ihre Geschwister verteilt hatten, und wer denn dieser Begleiter sei.

Gouri schüttelte Hände links und rechts und tauschte, inzwischen selbst in Tränen ausgebrochen, Umarmungen aus.

Dann bat sie um Ruhe, damit der edle ‚Prip, Sohn des Aigar, Bote des Eisernen Stuhls von Cavern und Hüter des Casamatucs der Drei Wasser‘ seinen Bericht der Ereignisse abgeben könne.

Prip, der sich während der ganzen Zärtlichkeiten peinlich berührt im Hintergrund gehalten hatte, trat nach vorne und räusperte sich. Er hatte auf dem Weg bereits daran gefeilt, wie er diese ganze Geschichte am besten vortragen könnte.

Der Alte Grone blieb am Rande der Menge stehen und gluckste fröhlich.

Die Troubadoure suchten bereits ihre Flöten hervor und putzten sie eifrig.

Die verlorene Tochter war zurückgekehrt.

Heute Abend würde es ein Fest geben.
